\documentclass[11pt,letterpaper]{article}
\usepackage[T1]{fontenc}
\usepackage{lmodern,microtype}
\usepackage[margin=1.04in]{geometry}
\usepackage{amsmath,amssymb,amsthm,mathtools}
\usepackage{enumitem}
\usepackage{listings}
\usepackage[hidelinks]{hyperref}
\lstset{basicstyle=\ttfamily\footnotesize,columns=fullflexible,
 keepspaces=true,showstringspaces=false,breaklines=true,
 literate={ℕ}{{\(\mathbb N\)}}1 {ℝ}{{\(\mathbb R\)}}1
 {ℂ}{{\(\mathbb C\)}}1 {→}{{\(\to\)}}2 {∀}{{\(\forall\)}}1
 {∃}{{\(\exists\)}}1 {∧}{{\(\land\)}}1 {≤}{{\(\le\)}}1
 {≠}{{\(\ne\)}}1 {∫}{{\(\int\)}}1 {∏}{{\(\prod\)}}1
 {∑}{{\(\sum\)}}1 {ξ}{{\(\xi\)}}1 {𝓕}{{\(\mathcal F\)}}1
 {‖}{{\(\Vert\)}}1 {⁻¹}{{\({}^{-1}\)}}1}
\hypersetup{pdftitle={Boundedness of higher-dimensional twisted paraproducts: a formalization blueprint}}
\newtheorem{theorem}{Theorem}[section]
\newtheorem{lemma}[theorem]{Lemma}
\newtheorem{proposition}[theorem]{Proposition}
\newtheorem{corollary}[theorem]{Corollary}
\theoremstyle{definition}
\newtheorem{definition}[theorem]{Definition}
\newcommand{\R}{\mathbb R}
\newcommand{\C}{\mathbb C}
\newcommand{\Z}{\mathbb Z}
\newcommand{\E}{\mathbb E}
\newcommand{\ind}{\mathbf1}
\newcommand{\dd}{\,\mathrm d}
\newcommand{\norm}[1]{\lVert#1\rVert}
\newcommand{\avg}[2]{\langle #1\rangle_{#2}}
\newcommand{\code}[1]{\texttt{\detokenize{#1}}}
\newcommand{\Sch}{\mathcal S}
\providecommand{\lean}[1]{}
\providecommand{\uses}[1]{}
\setlength{\emergencystretch}{1.6em}
\title{Boundedness of higher-dimensional twisted paraproducts:\\a formalization blueprint}
\author{}
\date{November 22, 2026}
\begin{document}
\maketitle

\section{The untruncated estimate}
Take \(d\ge2\), \(X=\R^d\) with Euclidean norm and Lebesgue measure,
and \(f_i\in\Sch(X;\C)\) for \(i\in[d]=\{1,\ldots,d\}\). Write
\(e_i\) for the \(i\)-th coordinate vector. For
\(\xi=(\xi^1,\ldots,\xi^d)\in X^d\), set
\(\eta(\xi)_i=\xi^i_i\). Use the Fourier character
\(e^{-2\pi i x\cdot\xi}\) for the forward transform. A smooth Mihlin
symbol on \(X\setminus\{0\}\) has finite seminorms
\[
 [m]_s=\max_{|\alpha|\le s}\sup_{\zeta\ne0}
            |\zeta|^{|\alpha|}\,|\partial^\alpha m(\zeta)|.
\]
Give \(m\) value zero at the origin. Define the untruncated operator by
\begin{equation}\label{bp:op}
 T_m(\vec f)(x)=\int_{X^d}m(\eta(\xi))
    \prod_{i=1}^d\widehat f_i(\xi^i)\,
    e^{2\pi i x\cdot\sum_i\xi^i}\dd\xi .
\end{equation}
For a smooth angular function \(\Omega\) of spherical mean zero put
\(K(y)=|y|^{-d}\Omega(y/|y|)\) for \(y\ne0\). The spatial version is
\[
 T_K(\vec f)(x)=\operatorname{p.v.}\int_X K(y)
                              \prod_i f_i(x-y_i e_i)\dd y .
\]
Here the principal value is the joint limit over
\(\epsilon<|y|<R\) as \(\epsilon\downarrow0,\ R\uparrow\infty\).

\begin{theorem}[A priori estimate on Schwartz functions]\label{bp:main}
\lean{TwistedParaproduct.proveBanachTarget}
\uses{bp:frechet,bp:full-multiplier,bp:homogeneous-symbol,bp:hard-homogeneous}
Suppose \(1<p_i<\infty\) and \(1/p=\sum_i1/p_i<1\).
There are \(s\in\mathbb N\) and \(C<\infty\), depending only on
\(d,\vec p\), such that for every \(m\) smooth away from the origin
with \([m]_s<\infty\), and all Schwartz inputs,
\[
                \norm{T_m(\vec f)}_{L^p(X)}
                \le C[m]_s\prod_{i=1}^d\norm{f_i}_{L^{p_i}(X)} .
\]
For every \(K\) above the untruncated principal value exists for these
inputs and equals \eqref{bp:op} with the principal-value Fourier symbol
of \(K\). In particular
\[
               \norm{T_K(\vec f)}_{L^p(X)}
                    \le C_{d,\vec p,\Omega}\prod_i\norm{f_i}_{L^{p_i}(X)} .
\]
\end{theorem}

The proof of Theorem~\ref{bp:main} is assembled after the identification
of the principal value. Appendix~\ref{bp:lean-target} gives the
typechecked Lean proposition for its Fourier estimate.
Appendix~\ref{sec:tpb-source-map} records exact library declarations
and the interfaces that require further formalization.
In the \LaTeX{} source, \verb|\lean| gives the proposed declaration name
and \verb|\uses| lists dependencies by label. The dyadic estimate is
Proposition~\ref{bp:dyadic}; the Fourier estimate and the pointwise
principal-value identification are Lemmas~\ref{bp:full-multiplier}
and~\ref{bp:hard-homogeneous}.

In Lean the
inputs can be a family of \code{SchwartzMap} indexed by \code{Fin d}.
Use \code{EuclideanSpace} for Fourier space and \code{eLpNorm} with
exponents \(\operatorname{ofReal}(p_i)\) and
\(\operatorname{ofReal}(p)\). Quantify a positive real constant and take
its image in \(\mathbb R_{\ge0}^{\infty}\) in the norm inequality.
All exponent reciprocals and all positive cutoffs in this text are
ordinary real numbers. For coordinate integration, the standard
orthonormal basis identifies \(X\) measure-preservingly with a finite
product of real lines. In that representation \(x[i:=y]\) replaces
coordinate \(i\). Evaluations of measurable functions below use
measurable representatives. The integration identities in the proof
descend to almost-everywhere equivalence classes.

\begin{lemma}[Well-defined Fourier model]\label{bp:schwartz-Fourier}
\lean{TwistedParaproduct.schwartzFourier}
For Schwartz inputs \eqref{bp:op} is absolutely convergent and
continuous in \(x\); its value is independent of the choice of \(m(0)\).
\end{lemma}
\begin{proof}
The Fourier transform of each \(f_i\) is Schwartz, hence integrable.
Fubini bounds the absolute integral by
\([m]_0\prod_i\norm{\widehat f_i}_1\). The set \(\eta=0\) is contained
in the coordinate hyperplane \(\xi^1_1=0\), which is null by Fubini.
The same integrable majorant gives continuity by dominated convergence.
\end{proof}

\begin{lemma}[Two equivalent symbol seminorms]\label{bp:frechet}
\lean{TwistedParaproduct.frechet}
Replacing every order-\(k\) coordinate derivative in \([m]_s\) by the
operator norm of the full \(k\)-linear Fr\'echet derivative gives an
equivalent seminorm, with constants depending only on \(d,s\).
\end{lemma}
\begin{proof}
Evaluate the Fr\'echet derivative on any ordered list of unit coordinate
vectors to obtain the corresponding iterated partial derivative.
Conversely, expand \(D^km(z)[v_1,\ldots,v_k]\) in the coordinate basis
on each of its \(k\) inputs. There are \(d^k\) terms, each bounded by
the largest coordinate derivative times \(\prod_j|v_j|\).
Smoothness on the open punctured space permits interchange of coordinate
derivatives. The \(k=0\) norm is \(|m(z)|\). Multiply by \(|z|^k\)
and take suprema. In particular, quantifying over all nonnegative
upper bounds \(B\) for the Fr\'echet version gives the same target.
\end{proof}

\section{Coordinate averages and finite martingales}

All dyadic intervals are half open. At integer level \(k\) they have
length \(2^{-k}\), and each interval is the disjoint union of its two
children. A translated nested grid works identically. Write
\(\langle f\rangle_I=|I|^{-1}\int_I f\) and
\(h_I=\ind_{I_{\rm left}}-\ind_{I_{\rm right}}\). The operator
\(E_k^i\) replaces a function by its average on the level-\(k\)
interval in the \(i\)-th coordinate, with the other coordinates fixed.
Write \(\Delta_k^i=E_{k+1}^i-E_k^i\).

\begin{lemma}[Conditional averages and Haar squares]\label{bp:haar-square}
\lean{TwistedParaproduct.haarSquare}
For \(1\le p<\infty\), \(E_k^i\) is an \(L^p\) contraction,
\(E_j^iE_k^i=E_{\min(j,k)}^i\), and on an interval \(I\) of
level \(k\) one has
\[
       \Delta_k^i f(x)=h_I(x_i)\langle h_I f\rangle_I .
\]
Here and below a one-coordinate average retains all the other arguments
of its integrand. For \(u<T\),
\(\sum_{k=u}^{T-1}\norm{\Delta_k^i f}_2^2\le\norm{f}_2^2\).
For any real \(C\in L^2(I)\) one also has
\[
 \tfrac12\sum_{J\ {\rm child\ of}\ I}\langle C\rangle_J^2
                  -\langle C\rangle_I^2=\langle h_I C\rangle_I^2 .
\]
\end{lemma}
\begin{proof}
Jensen's inequality on each interval, summation and Fubini give the
contraction. Nesting of intervals gives the composition rule. The two
child averages are \(\langle f\rangle_I\pm\langle h_I f\rangle_I\);
these two equalities give both Haar identities. Apply the square
identity to the real and imaginary parts of \(f\) and add.
Multiplying by \(|I|\), then summing and integrating
the other coordinates gives
\(\norm{\Delta_k^i f}_2^2=\norm{E_{k+1}^i f}_2^2-\norm{E_k^i f}_2^2\).
Sum in \(k\) and apply contraction.
\end{proof}

\begin{lemma}[Directional dyadic maximal inequality]\label{bp:dyadic-max}
\lean{TwistedParaproduct.dyadicMax}
Let \(M_i f=\sup_{k\in\mathbb Z}E_k^i|f|\). For \(f\in L^1(\R^d)\)
and \(\lambda>0\),
\[
          |\{M_i f>\lambda\}|\le \lambda^{-1}\norm{f}_1 .
\]
For almost every \(x\), \(E_k^i f(x)\to f(x)\) as \(k\to+\infty\).
\end{lemma}
\begin{proof}
On almost every fixed fiber, \(f\) is integrable. The averages of
\(|f|\) on intervals of lengths tending to infinity tend uniformly to
zero there. Thus the exceptional intervals with average exceeding
\(\lambda\) are contained in maximal exceptional intervals. These
are disjoint, and each has length at most \(\lambda^{-1}\) times its
integral of \(|f|\). Sum and integrate the fixed coordinates. All sets
are measurable since there are countably many dyadic intervals.
On the line the differentiation assertion holds for continuous functions
by continuity. For any \(g\in C_c(\R)\) the upper limit of
\(|E_k f-f|\) is bounded by \(M(f-g)+|f-g|\) almost everywhere.
The weak bound and Markov's inequality make the measure of this
upper limit exceeding \(\epsilon>0\) at most
\(4\epsilon^{-1}\norm{f-g}_1\). Approximation in \(L^1(\R)\)
and Fubini give the claim.
\end{proof}

\begin{lemma}[One-coordinate Calder\'on--Zygmund replacement]\label{bp:dyadic-cz}
\lean{TwistedParaproduct.dyadicCz}
Let \(0\le F\le\ind_A\), \(0<|A|<\infty\), and \(0<\lambda<1\).
On each \(i\)-fiber, replace \(F\) on the maximal intervals of
\(D=\{M_i\ind_A>\lambda\}\) by its average there, and keep its values
elsewhere. The resulting \(g\) satisfies, up to null sets,
\[
 0\le g\le 2\lambda,\qquad |\{g\ne0\}|\le |A|/\lambda,\qquad
 E_k^i g=E_k^i F,\quad\Delta_k^i g=\Delta_k^i F
                 \quad\text{on }D^c\quad(k\in\mathbb Z).
\]
\end{lemma}
\begin{proof}
\uses{bp:dyadic-max}
The exceptional intervals form the disjoint collection from
Lemma~\ref{bp:dyadic-max}. Their parents have density at most
\(\lambda\), so their own densities are at most \(2\lambda\).
Almost every point of \(A\) belongs to \(D\) by differentiation;
the measure estimate follows from the same lemma. Any dyadic
interval containing a point outside \(D\) either contains an entire
exceptional interval or is disjoint from it. The integral of \(F-g\)
on each exceptional interval vanishes. Absolute integrability
justifies summing these zero integrals. This applies at both levels
in the difference. Measurability follows by expressing maximality by
countably many comparisons on each fiber.
\end{proof}

We next give the finite martingale estimates used below. One may work
on a finite probability tree, since on a fixed initial dyadic interval
the process has finitely many terminal atoms. A finite martingale
\((M_k)_{k=0}^n\) on increasing partitions is constant on the atoms
at time \(k\); conditional averaging at time \(k\) returns \(M_k\)
from \(M_{k+1}\). All variables are real. Denote conditional
expectation by \(\E_k\), write \(D_k=M_{k+1}-M_k\) for \(0\le k<n\),
and put \(M^*=\max_{0\le k\le n}|M_k|\).

\begin{lemma}[Bounded optional sampling]\label{bp:optional}
\lean{TwistedParaproduct.optional}
For stopping times \(0\le\sigma\le\tau\le n\), \(M_\sigma\) is
measurable for the stopped partition and
\(\E[M_\tau\mid\mathcal F_\sigma]=M_\sigma\). In particular
\((M_{\tau_j})_j\) is a finite martingale for any finite
nondecreasing sequence of stopping times \(\tau_j\).
\end{lemma}
\begin{proof}
The stopped sigma algebra consists of events \(A\) such that
\(A\cap\{\sigma\le k\}\in\mathcal F_k\) for every \(k\).
Writing \(M_\sigma=\sum_k\ind_{\{\sigma=k\}}M_k\) proves the
measurability assertion. For such \(A\),
\[
 \E[\ind_A(M_\tau-M_\sigma)]
   =\sum_{k<n}\E[\ind_{A\cap\{\sigma\le k\}\cap\{\tau>k\}}D_k]=0 .
\]
The indicator in each term is measurable at time \(k\) and
\(\E_kD_k=0\). Applying the equality to successive times proves the
last assertion. This argument also proves
\(\E[\ind_A M_n]=\E[\ind_A M_\sigma]\) with \(\tau=n\).
\end{proof}

\begin{lemma}[Finite maximal inequality]\label{bp:doob}
\lean{TwistedParaproduct.doob}
For \(1<p<\infty\), \(\norm{M^*}_p\le p(p-1)^{-1}\norm{M_n}_p\).
\end{lemma}
\begin{proof}
\uses{bp:optional}
Stop at the first \(k\) with \(|M_k|>\lambda\), recording \(n\) if
there is no hit. The event \(A=\{M^*>\lambda\}\) is measurable at
this time. Lemma~\ref{bp:optional} and conditional Jensen imply
\(\lambda\mathbb P(A)\le\E[\ind_A|M_n|]\). Multiply by
\(p\lambda^{p-2}\), integrate in \(\lambda\), and apply Fubini
to get
\(\E(M^*)^p\le p(p-1)^{-1}\E[|M_n|(M^*)^{p-1}]\).
H\"older proves the assertion, including the zero-norm case.
\end{proof}

\begin{lemma}[Two finite square-function inequalities]\label{bp:mart-square}
\lean{TwistedParaproduct.martSquare}
Let \(S_k=(\sum_{j<k}|D_j|^2)^{1/2}\) and \(S=S_n\). For every
\(2\le p<\infty\),
\[
       \norm{S}_p\le C_p\norm{M_n}_p,\qquad
       \norm{M_n}_p\le C_p\norm{S}_p \quad\hbox{if }M_0=0 .
\]
\end{lemma}
\begin{proof}
\uses{bp:doob}
For \(p=2\) conditional averaging gives
\(\E S^2=\E M_n^2-\E M_0^2\). Take \(p>2\). For \(t,s\ge0\)
the mean value theorem, separately on \(s\le t\) and \(s>t\), gives
\((t+s)^{p/2}-t^{p/2}\le C_p(t^{p/2-1}s+s^{p/2})\).
Hence
\[
 \E S^p\le C_p \E\sum_{k<n}
       \big(S_k^{p-2}D_k^2+|D_k|^p\big).
\]
The expectation of the first sum equals
\(\E\sum_k S_k^{p-2}(M_{k+1}^2-M_k^2)\), since
\(\E_kD_k=0\). As \(S_k^{p-2}\) is increasing, discrete summation
by parts bounds this last sum pointwise by \(S^{p-2}(M^*)^2\).
The other sum is bounded by \((2M^*)^{p-2}S^2\). By H\"older,
if \(s=\norm{S}_p\), \(m=\norm{M^*}_p\), then
\(s^p\le C_p(s^{p-2}m^2+s^2m^{p-2})\). If \(s>m>0\), divide
by \(s^p\) and use the positive exponents \(2,p-2\) to get
\(s\lesssim_p m\). The other cases are immediate. Lemma~\ref{bp:doob}
proves the first estimate.

For the reverse estimate use the real-variable inequality
\[
 |a+b|^p-|a|^p-p|a|^{p-2}ab
       \le C_p\bigl(|a|^{p-2}|b|^2+|b|^p\bigr).
\]
It follows by integrating the second derivative of \(|\cdot|^p\)
between \(a\) and \(a+b\) and using
\(|a+tb|^{p-2}\lesssim_p |a|^{p-2}+|b|^{p-2}\).
Take \(a=M_k,b=D_k\), average, and telescope. The linear terms
integrate to zero, and \(\sum_k|D_k|^p\le S^p\). When \(M_0=0\),
\[
 \norm{M_n}_p^p\le C_p\big(\norm{M^*}_p^{p-2}\norm{S}_p^2+
                                     \norm{S}_p^p\big).
\]
Substitute Lemma~\ref{bp:doob} and divide by \(\norm{M_n}_p^p\)
when it is nonzero. The second assertion follows.
\end{proof}

\begin{lemma}[Coordinate martingale transforms]\label{bp:mart-transform}
\lean{TwistedParaproduct.martTransform}
For \(2\le p<\infty\), \(u<v\), and \(|c_k|\le1\) real, the
martingale \(Y_u=0\), \(Y_{k+1}=Y_k+c_k\Delta_k^b G\) satisfies
\[
                   \norm{Y_v}_p\le C_p\norm{G}_p .
\]
\end{lemma}
\begin{proof}
\uses{bp:mart-square,bp:haar-square}
On each normalized initial interval in coordinate \(b\), the increments
have conditional mean zero. For \(X_k=E_k^b G\) their square function
satisfies \(S(Y)\le S(X)\) pointwise. Both bounds of
Lemma~\ref{bp:mart-square} give
\(\norm{Y_v}_p\lesssim_p\norm{X_v}_p\). Raise the estimate to
the \(p\)-th power, multiply by the length of the initial interval,
sum over those intervals and integrate the other coordinates.
Apply Lemma~\ref{bp:haar-square} to \(X_v\). This argument also works
for coefficients measurable at time \(k\). Complex scalar coefficients
in a form can be treated by linearity in their real and imaginary parts.
\end{proof}

\subsection{Recording changes of amplitude}
Assume \(0\le G_i\le\ind_{A_i}\), \(a_i=|A_i|<\infty\).
Fix \(u<v\), let \(X^i_k=E_k^iG_i\) for
\(u\le k\le T=v+d+1\), \(i\ne b\), and let \(Y\) be the transform
of \(G_b\) on \(u\le k\le v\) from Lemma~\ref{bp:mart-transform}.
For a finite path \(M=(M_k)_{u\le k\le t}\) write
\(\operatorname{TV}(M)=\sum_{k=u}^{t-1}|M_{k+1}-M_k|\).
All scalar paths below are functions on \(\R^d\).

\begin{lemma}[Recording times, extinction, and variation]\label{bp:stopped-var}
\lean{TwistedParaproduct.stoppedVar}
Take \(M=X^i\) and \(a=a_i\), or \(M=Y\) and \(a=a_b\).
For each \(0<\alpha\le1\) there is a path \(Z^\alpha(M)\), adapted
to the coordinate filtration, with \(|M_k-Z^\alpha_k(M)|\le\alpha\) and
\begin{align*}
 \norm{\alpha\operatorname{TV}(Z^\alpha(M))}_q
      &\lesssim_q a^{1/q}\qquad(1\le q<\infty),\\
 \sum_k\int_{\R^d}|Z^\alpha_{k+1}(M)-Z^\alpha_k(M)|^2
      &\lesssim a .
\end{align*}
For \(M=X^i\) its values lie in \([0,1]\), and \(Z^\alpha_k(M)=0\)
whenever \(M_k=0\), up to null sets.
\end{lemma}
\begin{proof}
\uses{bp:optional,bp:mart-square,bp:mart-transform}
Start with \(Z_u=M_u\). At each later time record \(M_k\) if its
distance from the previous recorded value is at least \(\alpha\).
For \(X^i\), also record a new zero when the previous recorded value
was nonzero. Otherwise retain the previous value. These comparisons
use only values known by time \(k\), prove the tracking assertion,
and preserve \([0,1]\) for \(X^i\). A nonnegative martingale reaching
zero stays zero almost surely: on a zero atom, the average of its
nonnegative terminal variable is zero. Thus the extra recording of a
zero can happen at most once.

Number the recording times in increasing order, include the initial
time, and append copies of \(t\) so that the length is fixed and the
list includes the terminal time. A terminal time added by convention
does not change the recorded path. Each entry \(\tau_j\) is a
stopping time, since the list of hits up to time \(k<t\) is
measurable at \(k\). Set \(V=\sum_j |M_{\tau_{j+1}}-M_{\tau_j}|^2\).
Lemmas~\ref{bp:optional} and \ref{bp:mart-square} at exponent \(2q\)
give on each normalized initial interval
\(\norm{V}_q\lesssim_q\norm{M_t}_{2q}^2\). Raise to power \(q\),
sum with the interval lengths and integrate the fixed coordinates to
obtain the same inequality globally. Contraction for \(X^i\), or
Lemma~\ref{bp:mart-transform} for \(Y\), yields
\(\norm{M_t}_{2q}^{2}\lesssim_q a^{1/q}\).
Every jump of \(Z^\alpha\) is included in \(V\); \(q=1\)
gives the integrated square bound.

Every regular jump of size \(s\) satisfies \(\alpha s\le s^2\).
For the extra zero hit with jump size \(s<\alpha\), its cost
\(\alpha s\le\alpha^2\) is covered by an earlier regular jump
square if one exists. If none exists then \(s=|M_u|\). Consequently
\(\alpha\operatorname{TV}(Z^\alpha)\le2V+\alpha|M_u|\).
For the low path \(\norm{M_u}_q\le a^{1/q}\); the high path starts
at zero. This proves the variation bound.
\end{proof}

\begin{lemma}[Residual drift and Haar squares]\label{bp:residual-drift}
\lean{TwistedParaproduct.residualDrift}
For \(e_k=X^i_k-Z^\alpha_k(X^i)\), on each level-\(k\) interval
\(I\) of coordinate \(i\) there are uniquely determined functions
\(p_k,d_k\), constant on \(I\) with the other arguments fixed, with
\[
  e_{k+1}=e_k+p_k+h_I d_k,\qquad
  |p_k|\le2\alpha,\qquad |d_k|\le\alpha .
\]
They satisfy \(\sum_k\norm{p_k}_1\lesssim a_i/\alpha\) and
\(\sum_k\norm{d_k}_2^2\lesssim a_i\).
For the true path \(X^i\) itself the same child expansion has \(p_k=0\),
\(|d_k|\le1\), and \(\sum_k\norm{d_k}_2^2\le a_i\).
\end{lemma}
\begin{proof}
\uses{bp:stopped-var,bp:haar-square}
Take \(p_k=E_k^i e_{k+1}-e_k\) and let \(d_k\) be half the
difference of the two child values; the two-child identities give
uniqueness and both pointwise estimates. Martingality of \(X^i\)
gives \(p_k=-E_k^i(Z^\alpha_{k+1}-Z^\alpha_k)\). Apply contraction
and the \(q=1\) variation bound of Lemma~\ref{bp:stopped-var}.
Also \(|d_k|^2\le E_k^i|e_{k+1}-e_k|^2\) by Jensen. The latter
averages are bounded by twice the sum of the squared increments of
\(X^i\) and \(Z^\alpha\), conditionally averaged. Integrate and sum
using Lemmas~\ref{bp:haar-square} and \ref{bp:stopped-var}. The
assertions for \(X^i\) follow directly from Lemma~\ref{bp:haar-square}.
\end{proof}

\begin{lemma}[All recorded low paths against the high transform]\label{bp:all-coarse}
\lean{TwistedParaproduct.allCoarse}
Suppose \(d\ge2\), \(0\le H\le\ind_{A_0}\), \(h=|A_0|>0\),
\(a_i\ge h\), and \(R=\max_{i\in[d]}a_i/h\). For an integer \(N\ge0\)
put \(\delta=2^{-N}\). For \(1<q<\infty\),
\[
 \left|\int_{\R^d}H(x)\sum_{k=u}^{v-1}
       \left(\prod_{i\ne b} Z^\delta_k(X^i)(x)\right)
                           (Y_{k+1}(x)-Y_k(x))\,dx\right|
          \lesssim_{d,q}(N+1)^d h R^{1/q}.
\]
\end{lemma}
\begin{proof}
\uses{bp:stopped-var}
For each path \(M\) write \(M^{[0]}=Z^1(M)\) and
\(M^{[j]}=Z^{2^{-j}}(M)-Z^{2^{1-j}}(M)\) for \(1\le j\le N\).
These paths add to \(Z^\delta(M)\). The tracking inequality gives
\(\sup_k|M_k^{[j]}|\le3\,2^{-j}\) for \(j>0\). For low paths
this estimate also holds for \(j=0\), since \(0\le X^i\le1\).
For all \(j\), Lemma~\ref{bp:stopped-var} gives
\(\norm{\operatorname{TV}(M^{[j]})}_q\lesssim_q 2^j a^{1/q}\).
H\"older therefore gives
\[
        \int H\operatorname{TV}(M^{[j]})
              \lesssim_q 2^j h^{1-1/q}a^{1/q}
              \le C_q\,2^j hR^{1/q}.
\]

First take \(Z^\delta(Y)\) in the high position. Choose low levels
\(n_i\) and high level \(m\), all in \(\{0,\ldots,N\}\), and put
\(L_k=\prod_{i\ne b}(X^i)^{[n_i]}_k\). If \(m\le\max_{i\ne b}n_i\),
the absolute pointwise sum is bounded by
\(\sup|L|\,\operatorname{TV}(Y^{[m]})\). The integral costs at most
\(C_{d,q}2^{m-\sum_{i\ne b}n_i}hR^{1/q}\le C_{d,q}hR^{1/q}\).
This uses only variation for the possibly unbounded \(Y^{[0]}\).

If \(m>\max_{i\ne b}n_i\), then \(m\ge1\). The identity obtained
by expanding \(L_{k+1}W_{k+1}-L_kW_k\) and summing yields
\[
 \left|\sum_{k=u}^{v-1}L_k(W_{k+1}-W_k)\right|
       \le \sup_k|W_k|\bigl(2\sup_k|L_k|+\operatorname{TV}(L)\bigr).
\]
Apply this with \(W=Y^{[m]}\). Replacing one product factor at a
time gives
\[
 \operatorname{TV}(L)\le
  \sum_{r\ne b}\operatorname{TV}((X^r)^{[n_r]})
                \prod_{i\ne b,r}\sup_k|(X^i)_k^{[n_i]}|.
\]
The endpoint contribution is \(O_d(h)\) after integration. The
\(r\)-th variation contribution costs at most
\(C_{d,q}2^{-m+n_r-\sum_{i\ne b,r}n_i}hR^{1/q}
\le C_{d,q}hR^{1/q}\). Summing all \((N+1)^d\) tuples gives
the desired bound for \(Z^\delta(Y)\).

Finally apply the same summation by parts with \(W=Y-Z^\delta(Y)\),
so \(\sup|W|\le\delta\), and \(L_k=\prod_{i\ne b}Z^\delta_k(X^i)\).
The low factors lie in \([0,1]\), whence
\(\operatorname{TV}(L)\le\sum_{i\ne b}\operatorname{TV}(Z^\delta(X^i))\).
H\"older and Lemma~\ref{bp:stopped-var} bound the integral by
\(C_{d,q}hR^{1/q}\).
\end{proof}

\paragraph{Implementation.}
Use one finite tree at a time for Lemmas~\ref{bp:optional}--\ref{bp:mart-square}.
Conditional expectations on a tree are weighted sums on atoms.
Represent a recording time by a natural number bounded by the terminal
index. The recorded path can be defined by updating the last value at
each index; the successive stopping times can be obtained by counting
the updates. Extend their list by the terminal index, as in
Lemma~\ref{bp:stopped-var}. For an interval offset use the low path at
the actual generation of that interval. All \(L^p\) statements on
\(\R^d\) follow by integration of finite-tree inequalities; this avoids
putting a probability normalization on \(\R^d\) itself.

\section{Pruning and reflected rectangles}
Fix \(u<v\) in \(\Z\), \(b\in[d]\), and scalars \(|c_k|\le1\). For now
use real coefficients and real, nonnegative functions
\(0\le F_i\le\ind_{A_i}\) for \(0\le i\le d\), with finite positive
measures \(a_i=|A_i|\). Write \(H=F_0,\ h=a_0\), and
\[
  \Lambda(\vec F)=\sum_{k=u}^{v-1}c_k
       \int H\,\Delta_k^b F_b\prod_{i\in[d]\setminus\{b\}}E_k^iF_i .
\]
Fix \(t_0,\ldots,t_d>0\) of sum one and put \(W=\prod_{i=0}^d a_i^{t_i}\).
Define \(C_*\) as the supremum of \(|\Lambda|/W\) over these data
and coefficients with \(b,u,v,\vec t\) fixed. This supremum exists as
a finite real number: coordinate averages are \(L^q\)-contractions by
Jensen and Fubini, so H\"older with \(q_i=1/t_i\) gives
\(C_*\le2(v-u)\). Terms involving null supports are zero by the same
inequality.

\begin{lemma}[Pruning and monotone masks]\label{bp:prune}
\lean{TwistedParaproduct.prune}
Assume \(h\le a_i\) for \(i\in[d]\). For \(K_*>1\) one can prune
each \(F_i\) to \(G_i\) at total form cost
\(dK_*^{-\min_{i\ge1}t_i}C_*W\). On each active dyadic cube \(Q\)
there is a tester \(0\le H_Q\le H\) with the following properties.
Its summed parent-to-child mass loss, including the final drop to
zero, is at most \(h\). In coordinate \(i\), on a subinterval of the
\(i\)-side of \(Q\) of relative length at least \(2^{-(d+1)}\),
its fiber average is at most \(2^{d+1}K_*h/a_i\). In a pruned summand
one can replace \(H\) by \(H_Q\). This remains true on replacing each
low average by either itself or its zero-reset residual.
\end{lemma}
\begin{proof}
\uses{bp:dyadic-max,bp:stopped-var,bp:haar-square}
Put \(U_i=\{M_i\ind_{A_0}>K_*h/a_i\}\) for the directional dyadic
maximal function and \(G_i=F_i\ind_{U_i^c}\).
Lemma~\ref{bp:dyadic-max} gives \(|U_i|\le a_i/K_*\).
Change one slot at a time. The changed difference has support
\(A_i\cap U_i\), giving the asserted cost by the definition of \(C_*\).

At a point of \(Q\), keep \(H(x)\) precisely if, for every ancestor
\(Q'\supseteq Q\) including \(Q\), and every \(s\in[d]\),
\[
       \avg{\ind_{A_0}(x[s:=\,\cdot\,])}{I_s(Q')}\le K_*h/a_s .
\]
These countable tests are measurable and only accumulate on descending
chains. Give every terminal cube the zero tester. Pointwise telescoping
from \(u\) to \(v\), followed by integration, proves the loss bound.
The test on \(Q\) in coordinate \(s\) is constant along that fiber.
If it passes, the integral on a subinterval is at most the bound for
the whole interval times the ratio of lengths after normalization.
If it fails, the mask is zero there. If any ancestor test fails, all
points of the original \(s\)-interval lie in \(U_s\) on this fiber.
Consequently both child averages of \(G_s\) vanish when \(s=b\);
when \(s\ne b\) the parent average vanishes. The reset convention
also kills its residual there. This proves all the replacement
identities.
\end{proof}

Fix \(\delta=2^{-N}\le1/2\) and a nonempty set
\(D\subseteq[d]\setminus\{b\}\). On low slot \(i\) take
\(L_i=e_i^\delta\) for \(i\in D\) and \(L_i=X_i=E^iG_i\) for the
others. Write \(\sigma_i=\delta\) or \(1\), respectively, and use
the actual interval value \(L_{i,I}(x_{\hat i})\).
The two child values can be written \(L_{i,I}+P_{i,I}\pm D_{i,I}\),
where \(P_{i,I}\) is the average increment and \(D_{i,I}\) the
half-difference. The symbol \(D_{i,I}\) denotes a function here.
By Lemmas~\ref{bp:haar-square} and~\ref{bp:residual-drift},
\[
 |L_{i,I}|,|D_{i,I}|\le\sigma_i,\quad |P_{i,I}|\le2\sigma_i,\quad
 \sum_k\int |D_{i,k}|^2\lesssim a_i,\quad
 \sum_k\int |P_{i,k}|\lesssim a_i/\delta\quad(i\in D).
\]
On a true martingale \(P_{i,I}=0\). Use path values from \(u\) through
\(v+d+1\) in these estimates.

\begin{definition}[Corner measure and marked energy]\label{bp:energy-def}
\lean{TwistedParaproduct.energyDef}
\uses{bp:prune,bp:residual-drift}
For a rectangle \(V=\prod I_i\) use the normalized Lebesgue measure
on each interval. For \(b\in J\subseteq[d]\) give coordinates
\(j\in J\) two independent samples \(x_j^0,x_j^1\), and coordinates
outside \(J\) one sample each. Use the product probability measure.
For \(\omega:J\to\{0,1\}\) select the corner \(x^\omega\) by these bits.
Write \(n=2^{|J|}\) and
\[
 B_J(V;Q)=\E\prod_\omega
        \left[H_Q(x^\omega)\prod_{i\notin J}L_{i,I_i}(x^\omega_{\hat i})\right].
\]
Let \(S_i\) average over the two \(i\)-children of \(V\), evaluating
the box with the same mask \(H_Q\); set \(\partial_i=S_i-1\).
For offsets \(\nu_i\in\{0,\ldots,|J|-1\}\), write \(\mathcal V_{k,\nu}\)
for all dyadic rectangles with lengths \(2^{-k-\nu_i}\). Each has a
unique enclosing level-\(k\) cube \(Q_k(V)\). Define
\[
 U(J)=\max_{\nu,\,j\in J}
 \sum_{k=u}^{v-1}\sum_{V\in\mathcal V_{k,\nu}}
        |V|\,\partial_jB_J(V;Q_k(V)).
\]
Each grid is countable. In Lean one can use one sample per member of the
disjoint sum of \(\operatorname{Fin}d\) and \(J\), adding a second
sample only in \(J\). Corner selection is a measurable map out of this
product.
\end{definition}

\begin{lemma}[Marginals and countable grid sums]\label{bp:marginal}
\lean{TwistedParaproduct.marginal}
Each corner has normalized Lebesgue measure on \(V\) as its marginal
law. Each fixed-level grid sum in the energy is absolutely convergent.
\end{lemma}
\begin{proof}
\uses{bp:energy-def}
For every bounded measurable function of that corner, integrate out
the unselected coordinates by Fubini. Their probabilities have mass
one. The same identity holds for nonnegative measurable functions by
monotone convergence. Bound every factor of a box by one except one
central factor, which is at most \(\ind_{A_0}\). Multiply its integral
by the rectangle volume and sum over the grid to obtain at most \(h\).
The same bound holds for child averages since child volume is parent
volume times the averaging weight. The triangle inequality for their
difference proves the second assertion.
\end{proof}

\begin{lemma}[One reflected coordinate]\label{bp:box-square}
\lean{TwistedParaproduct.boxSquare}
The box and \(\partial_jB_J\) are nonnegative for \(j\in J\).
The pruned form with tester \(H_Q\), high input \(G_b\), and low
paths \(L_i\) has absolute value at most \(\sqrt{a_bU(\{b\})}\).
\end{lemma}
\begin{proof}
\uses{bp:haar-square,bp:energy-def,bp:marginal}
Hold all samples except the two \(j\)-samples fixed. The corner product
splits into two copies of the same real function \(C(y)\), so its
conditional integral is \(\avg C{I_j}^2\). Write
\(h_I=\ind_{I_{\rm left}}-\ind_{I_{\rm right}}\). For child means
\(a,b'\), the mean of the two squares minus the parent mean squared
is \((a-b')^2/4=\avg{h_IC}{I}^2\). This proves positivity for signed
low paths. On the unstaggered grid the local high difference is
\(h_{I_b}(x_b)\avg{h_{I_b}G_b}{I_b}\) on each fiber.
Integrate in \(x_b\) and apply Cauchy--Schwarz over the transverse
points and marked cubes with weight \(|Q|\) using normalized fiber
measures. The first square sums to at most \(a_b\) by
Lemma~\ref{bp:haar-square}. The second sums to
\(\sum_{k,Q}|Q|\partial_bB_{\{b\}}(Q;Q)\) by the identity just proved.
\end{proof}

\begin{lemma}[Telescoping moving masks]\label{bp:box-boundary}
\lean{TwistedParaproduct.boxBoundary}
At fixed \(J,\nu\) the grid sum with weights \(|V|\) of
\((\prod_{i=1}^d S_i-1)B_J(V;Q_k(V))\) is at most
\(nh\prod_{i\notin J}\sigma_i^n\). A marked derivative is bounded
above by this quantity plus the absolute contributions from
unreflected-coordinate refinements on grids with offsets at most
\(|J|\).
\end{lemma}
\begin{proof}
\uses{bp:prune,bp:box-square,bp:marginal}
Refining once in every direction sends the grid at \(k\) exactly to
the one at \(k+1\), with each child weight its volume. Each child
rectangle lies in a unique level-\((k+1)\) cube whose parent is the
old enclosing cube. Insert its new mask and telescope. All low factors
are functions of their actual intervals and require no replacement.
The initial boxes enter with minus sign and the last new boxes vanish.
In a remaining difference, change the \(n\) masks one at a time. Keep
only the changing mask difference, bound the other central factors by
one and retain all low amplitude bounds. The integrated difference at
one corner is the normalized integral of that mask drop on the
rectangle. The child rectangles partition the new cubes.
Lemma~\ref{bp:prune} bounds the total integrated drop by \(h\).
This proves the first assertion.

The operators \(S_i\) in distinct coordinates commute by finite
averaging. Expand their product minus the identity in telescoping
order, starting with the marked \(j\). Each other reflected term,
preceded by positive averages, is nonnegative by
Lemma~\ref{bp:box-square}. Move the unreflected terms to the other side.
Their preceding averages refine selected coordinates once. This
raises their offsets by at most one and preserves the original
level-\(k\) mask.
\end{proof}

\begin{lemma}[Unreflected-coordinate expansion]\label{bp:box-expand}
\lean{TwistedParaproduct.boxExpand}
Fix \(i\notin J\) and one of the staggered grids produced in
Lemma~\ref{bp:box-boundary}. Up to constants depending on \(d\), an
individual term in the sum of \(\partial_i B_J\) with a drift or at
least two half-differences costs at most
\[
    a_i\sigma_i^{n-2}\prod_{\ell\notin J,\ \ell\ne i}\sigma_\ell^n.
\]
A single half-difference and no drift costs at most
\(\sigma_i^{n-1}\sqrt{a_i U(J\cup\{i\})}\).
\end{lemma}
\begin{proof}
\uses{bp:energy-def,bp:marginal,bp:residual-drift,bp:box-square,bp:box-boundary}
For each corner write \(A_\omega,P_\omega,D_\omega\) for its low
value of type \(i\), drift and half-difference; these depend only on
the transverse samples. Write \(C(y)\) for all remaining factors,
with the common base \(i\)-coordinate set to \(y\). At a child of
sign \(\tau\) expand \(\prod_\omega(A_\omega+P_\omega+\tau D_\omega)\).
Averaging over the two children integrates \(C(y)\) against
\(h_{I_i}(y)^m\), where \(m\) is the number of selected \(D\)'s.
The all-\(A\) term cancels on subtracting the parent.

In a term with a drift, retain one \(|P_\alpha|\) and bound the other
\(n-1\) factors of type \(i\) by \(2\delta\). After discarding the
central factors, the retained corner marginal is uniform on its
rectangle. At fixed offset the \(i\)-levels \(k+\nu_i\) are distinct.
Summing its integral therefore costs at most \(C a_i/\delta\).
For a term with no drift and at least two \(D\)'s retain two and use
\(2|D_\alpha D_\beta|\le D_\alpha^2+D_\beta^2\). The grid sum for
either square is \(O(a_i)\) by the same marginal calculation.
Every other low type still occurs \(n\) times. This proves the first
bound.

For a single \(D_\alpha\), apply Cauchy--Schwarz to the sum over
rectangles and transverse samples. Put \(D_\alpha\prod_{\omega\ne\alpha}
A_\omega\) in the first factor; it costs
\(O(\sigma_i^{n-1}\sqrt{a_i})\). The second square integrates
\(\avg{h_{I_i}C}{I_i}^2\). Identifying the two copies of its integration
variable with the two independent \(i\)-samples makes it exactly a
marked \(i\)-derivative in \(J\cup\{i\}\). All its offsets are at most
\(|J|\), as allowed in that energy.
\end{proof}

\begin{lemma}[The two-copy density gain]\label{bp:two-copy}
\lean{TwistedParaproduct.twoCopy}
If \(J=\{b\}\) and \(i\in D\), each drift or multiple-difference term
of Lemma~\ref{bp:box-expand} has the sharper bound
\(C_{d,K_*}h^2/a_b\).
\end{lemma}
\begin{proof}
\uses{bp:prune,bp:box-expand,bp:marginal,bp:residual-drift}
Retain both central factors. After the drift or square estimate, the
remaining coefficient at one anchor \(\alpha\) is bounded by a
constant times \(\delta |P_\alpha|\) or \(D_\alpha^2\).
It depends on the anchor's \(b\)-sample and the other transverse
coordinates, and uses neither the opposite \(b\)-sample nor the base
\(i\)-sample. First fix the base \(i\)-sample \(y\). Integrate the
opposite central factor in its independent \(b\)-sample at cost
\(2^{d+1}K_*h/a_b\). Integrate the anchor central factor in \(y\) at
cost \(2^{d+1}K_*h/a_i\). These are precisely the uniform subinterval
bounds in Lemma~\ref{bp:prune}; they apply with the inherited mask.
The retained coefficient now sums to \(O(a_i)\) as in
Lemma~\ref{bp:box-expand}. Multiplication yields the claim.
\end{proof}

\begin{lemma}[Residual estimate]\label{bp:residual}
\lean{TwistedParaproduct.residual}
Set \(M=\max_i a_i=hR\). For every nonempty choice \(D\) the residual
form of Lemma~\ref{bp:box-square} is bounded in absolute value by
\(C_{d,K_*}(h+M\delta^{2^{-d}})\).
\end{lemma}
\begin{proof}
\uses{bp:box-boundary,bp:box-expand,bp:two-copy,bp:box-square,bp:energy-def}
Induct on the number of unreflected coordinates. By
Lemmas~\ref{bp:box-boundary} and~\ref{bp:box-expand}, \(U(J)\lesssim_d M\):
the terminal boundary costs \(O_d(h)\), each multiple term costs
\(O_d(a_i)\) and each singleton costs \(O_d(\sqrt{a_iM})\).
The number of expansion terms is at most a constant depending on \(d\).

For \(|J|\ge2\) and \(D\setminus J\ne\varnothing\), a second induction
gives \(U(J)\lesssim_d M\delta^{2^{-(d-|J|)}}\).
The boundary has a factor \(\delta^n\). Multiple terms have at least
\(\delta^{n-2}\) since some residual occurs and \(n\ge4\).
A singleton at a residual type has the gain \(\delta^{n-1}\) even with
the crude child bound. A singleton at an ordinary type keeps a residual
in the child and halves its exponent by the square root.
The assertion at \(J=[d]\) has empty premise.
At \(J=\{b\}\) the boundary and ordinary multiple terms have a factor
\(\delta^2\). A residual multiple term costs \(O_{d,K_*}(h^2/a_b)\)
by Lemma~\ref{bp:two-copy}. A residual singleton gains \(\delta\).
An ordinary singleton inherits the square root of the strong child
bound. If \(d=2\) such an ordinary singleton cannot occur. Thus
\[
 U(\{b\})\lesssim_{d,K_*} h^2/a_b+M\delta^{2^{-(d-1)}} .
\]
Lemma~\ref{bp:box-square}, \(a_b\le M\), and subadditivity of the square
root give the asserted estimate.
\end{proof}

\section{Uniform strong estimates for the dyadic form}
\begin{lemma}[Enlarging a supporting set]\label{bp:pad}
\lean{TwistedParaproduct.pad}
If \(A\subset X\) is measurable and \(|A|\le h<\infty\), then there
is a measurable \(B\supseteq A\) with \(|B|=h\).
\end{lemma}
\begin{proof}
The function \(r\mapsto|[-r,r]^d\setminus A|\) is continuous on
\([0,\infty)\) by dominated convergence on bounded ranges of \(r\);
cube boundaries have volume zero by Fubini. It starts at zero and
tends to infinity since \(|A|<\infty\). By the intermediate value
theorem choose \(r\) giving \(h-|A|\) and adjoin that cube to \(A\).
\end{proof}

\begin{lemma}[Uniform restricted form estimate]\label{bp:restricted}
\lean{TwistedParaproduct.restricted}
For the data defining \(C_*\), \(C_*\lesssim_{d,\vec t}1\) uniformly
in \(b,u,v\). This restricted bound also holds for complex coefficients
and complex inputs satisfying \(|F_i|\le\ind_{A_i}\).
\end{lemma}
\begin{proof}
\uses{bp:prune,bp:all-coarse,bp:residual,bp:dyadic-max,bp:dyadic-cz,bp:pad,bp:haar-square,bp:stopped-var}
First assume \(h\le\min_{i\ge1}a_i\). Apply Lemma~\ref{bp:prune} and
use \(R=\max_{i\ge1}(a_i/h)\). Take
\[
 N=\max(1,\lceil 2^d\log_2 R\rceil),\qquad \delta=2^{-N}.
\]
For \(I=[d]\setminus\{b\}\), the finite product identity
\[
 \prod_{i\in I}X_i-\prod_{i\in I}Z_i^\delta
   =\sum_{\emptyset\ne D\subseteq I}(-1)^{|D|+1}
           \prod_{i\in D} e_i^\delta \prod_{i\in I\setminus D}X_i
\]
follows by expanding \(\prod(X_i-e_i^\delta)\). Keep the tester \(H\)
in the all-recorded term of Lemma~\ref{bp:all-coarse}. Insert \(H_Q\)
in every other term by Lemma~\ref{bp:prune}, and apply
Lemma~\ref{bp:residual}. Their total cost is \(O_{d,K_*}(h)\) since
\(R\delta^{2^{-d}}\le1\).
Set \(t_*=\min_{i\ge1}t_i\). For \(q>2/t_*\), the coarse cost is bounded by
\[
 C_{d,q}(1+N)^d hR^{1/q}
     \lesssim_{d,\vec t} hR^{t_*}
     \le h\prod_{i=1}^d(a_i/h)^{t_i}=W .
\]
The middle inequality uses boundedness of
\((1+x)^d e^{-(t_*-1/q)x}\) on \(x\ge0\); the last uses \(a_i/h\ge1\).
Adding back the discarded data bounds \(|\Lambda|/W\) by
\(C_{d,\vec t,K_*}+dK_*^{-t_*}C_*\) in this case.

For arbitrary sizes choose \(L>d\). For each \(i\) with \(a_i<h/L\)
remove \(V_i=\{M_i\ind_{A_i}>L a_i/h\}\) from the tester.
The union has measure at most \(dh/L\) by Lemma~\ref{bp:dyadic-max}.
The tester on this union costs at most \((d/L)^{t_0}C_*W\).
Outside the union use Lemma~\ref{bp:dyadic-cz} on each such input.
It preserves all factors at every \(x\) in the remaining tester,
including the difference if \(i=b\). The replacement has amplitude at
most \(2L a_i/h\) and support of measure at most \(h/L\).
Divide it by \(2L a_i/h\) and enlarge its supporting set to size \(h\)
by Lemma~\ref{bp:pad}. For \(h/L\le a_i<h\) just enlarge the supporting
set to size \(h\). Keep \(A_0\) as a supporting set for the remaining
tester. This permits the first case. Restoring the normalization costs,
relative to \(W\), at most \(2L(a_i/h)^{1-t_i}\le2L\) per
replaced slot and \((h/a_i)^{t_i}\le L\) per enlarged slot. Hence
\[
 C_* \le (2L)^d C_{d,\vec t,K_*}
        +\big[(d/L)^{t_0}+(2L)^d d K_*^{-t_*}\big] C_* .
\]
Choose \(L\) and then \(K_*\) so the bracket is at most \(1/2\).
The initial \(C_*<\infty\) allows subtraction. Finally expand each
complex function into positive and negative real and imaginary parts
and split the sequence of complex coefficients into real and imaginary
parts. Every real part so obtained is covered by the restricted
estimate after taking out its sign. This costs a constant depending
only on the number of slots.
\end{proof}

\begin{lemma}[Restricted interpolation in the open H\"older simplex]
\label{bp:interp}
\lean{TwistedParaproduct.interp}
Let an \(m\)-linear scalar form admit restricted estimates on
\(|g_i|\le\ind_{B_i}\) with bound \(C_{\vec s}\prod_i |B_i|^{s_i}\)
at a relative neighborhood of a positive tuple
\(\vec t=(t_1,\ldots,t_m)\) of sum one. Then on simple functions of
finite-measure support the form is bounded by
\(C\prod_i\norm{g_i}_{1/t_i}\).
\end{lemma}
\begin{proof}
Fix \(\rho>0\) small enough that each
\(\vec t+\rho(e_i-e_j)\) lies in the neighborhood and use the largest
of these finitely many constants and that at \(\vec t\).
Split \(g_i\) into the disjoint amplitude layers \(g_{i,j}\) on
\(B_{i,j}=\{2^j\le |g_i|<2^{j+1}\}\). Ignore null layers and set
\(k_{i,j}=\lfloor\log_2 |B_{i,j}|\rfloor\). These lists are finite
for simple functions. Define
\[
  b_i(k)=2^{t_i k}\sum_{\{j:\,k_{i,j}=k\}}2^j.
\]
In each nonempty group the sum of the distinct powers \(2^j\)
is at most twice the largest. Therefore
\(\sum_k b_i(k)^{1/t_i}\lesssim_{\vec t}
\sum_j2^{j/t_i}|B_{i,j}|\le\norm{g_i}_{1/t_i}^{1/t_i}\).
At a fixed tuple of support bins choose index \(r\) of maximal \(k_i\)
and \(s\) of minimal \(k_i\). Apply the restricted estimate at
\(\vec t+\rho(e_s-e_r)\) to the normalized layers.
Since \(2^{k_i}\le |B_{i,j_i}|<2^{k_i+1}\), summing their amplitudes
within each bin costs at most a universal factor times
\[
              2^{-\rho(\max k_i-\min k_i)}\prod_i b_i(k_i).
\]
For offsets \(l_i=k_i-k_1\) with \(l_1=0\), H\"older in \(k_1\) bounds
the sum of products by \(\prod_i\norm{b_i}_{\ell^{1/t_i}}\).
The number of offset vectors with \(\max_i|l_i|\le n\) is at most
\((2n+1)^{m-1}\), and \(\max k_i-\min k_i\ge\max_i|l_i|\).
The resulting polynomial times geometric series converges. This proves
the lemma with no weak-space endpoint.
\end{proof}

\begin{proposition}[Strong dyadic bound]\label{bp:dyadic}
\lean{TwistedParaproduct.dyadic}
For \(1<p_i<\infty\), \(\sum_{i=0}^d1/p_i=1\),
\[
 \left|\sum_{k=u}^{v-1}c_k\int F_0\,\Delta_k^b F_b
                     \prod_{i\in[d]\setminus\{b\}}E_k^i F_i\right|
          \lesssim_{d,\vec p} \prod_{i=0}^d\norm{F_i}_{p_i}.
\]
This estimate holds for \(L^{p_i}\) inputs, all finite scale intervals,
and \(|c_k|\le1\). It has the same constant after translating the
grids independently and dilating them by a fixed positive factor in
each coordinate.
\end{proposition}
\begin{proof}
\uses{bp:restricted,bp:interp,bp:haar-square}
Use Lemmas~\ref{bp:restricted} and~\ref{bp:interp}
with \(m=d+1\). Approximate each input in \(L^{p_i}\) by simple
functions with finite-measure support. For fixed endpoints the
elementary \(2(v-u)\) bound controls the sum of the telescoping
approximation errors. Passing to the limit retains the uniform
constant. Conjugate by the diagonal affine map
\(x_i\mapsto a_i+\lambda_i x_i\) for translated and dilated grids.
The form picks up one Jacobian \(\prod_i\lambda_i\), and the norm
product the same factor since \(\sum_i1/p_i=1\).
\end{proof}

\section{From coordinate martingales to Fourier multipliers}
\label{bp:fourier-section}

Use the convention \(\widehat g(\xi)=\int g(x)e^{-2\pi i x\cdot\xi}\dd x\).
Fix \(d\ge2,\ 1<p_i<\infty\), and \(p>1\) with
\(1/p=\sum_{i=1}^d1/p_i\). All input functions in the final multiplier
statement are Schwartz functions on \(\R^d\). We use the elementary
Fourier inversion and Plancherel theorems, Fubini--Tonelli,
Lebesgue differentiation, and the usual integral inequalities for
Lebesgue spaces. The estimates below give the other one-dimensional
analysis needed in the reduction.

\subsection{Finite sums of one-dimensional kernels}
For a locally absolutely continuous \(a:\R\to\mathbb C\), define
\[
 Q(a)=\mathop{\rm ess\,sup}_{y\in\R}
 (1+|y|)^4\bigl(|a(y)|+|a'(y)|\bigr),\qquad
 a_k(y)=2^k a(2^k y)\quad(k\in\mathbb Z).
\]
Write \(a*_i f(x)=\int_\R a(y)f(x-ye_i)\dd y\) and
\(A(y)=(1-|y|)_+\). In particular \(Q(A)<\infty\) and \(\int A=1\).

\begin{lemma}[Maximal kernels on a coordinate]\label{bp:profile-max}
\lean{TwistedParaproduct.profileMax}
For \(1<q<\infty\),
\[
 \left\|\sup_{k\in\mathbb Z}|a_k*_i f|\right\|_q
 \le C_q Q(a)\|f\|_q .
\]
\end{lemma}
\begin{proof}
On the line let \(M g(x)=\sup_{s>0}(2s)^{-1}
\int_{x-s}^{x+s}|g|\). Integrating first on \(|y|\le 2^{-k}\)
and then on \(2^{j-k}<|y|\le 2^{j+1-k}\), \(j\ge0\), bounds
\(|a_k*g|\) by \(C Q(a)(1+\sum_{j\ge0}2^{-3j})Mg\).
The maximal function obeys
\(|\{Mg>\lambda\}|\le 3\|g\|_1/\lambda\).
Indeed each point of this open set admits a centered interval with
average greater than \(\lambda\). Cover a compact subset by finitely
many such intervals and select intervals of largest remaining length,
discarding those that intersect a selected one. The threefold dilates
of the disjoint selected intervals cover the compact set. Inner
regularity gives the claim. For \(g\in L^q\), the part of \(g\) of
modulus at most \(\lambda/2\) has maximal function at most \(\lambda/2\).
Apply the weak estimate to the rest and integrate the distribution
bound against \(q\lambda^{q-1}\dd\lambda\). Tonelli and
\(\int_0^{2|g|}\lambda^{q-2}\dd\lambda=(2|g|)^{q-1}/(q-1)\)
give the strong bound. Applying this argument on lines and integrating
in the spectator coordinates proves the assertion. Suprema may be
taken over rational radii by continuity in the radius.
\end{proof}

\begin{lemma}[Scalar convolution bound from cancellation]
\label{bp:scalar-cz}
\lean{TwistedParaproduct.scalarCz}
Suppose \(K\in L^1(\R)\), \(\|\widehat K\|_\infty\le B\), and
\[
 |K(x-u)-K(x)|\le B|u|\,|x|^{-2}\qquad(2|u|\le |x|,\ x\ne0).
\]
For \(1<q<\infty\), \(\|K*g\|_q\le C_q B\|g\|_q\),
independently of \(\|K\|_1\).
\end{lemma}
\begin{proof}
\uses{bp:dyadic-max}
Normalize \(B=1\). The case \(B=0\) follows by Fourier uniqueness.
Plancherel gives the \(L^2\) bound. For \(g\in L^1\) take the maximal
dyadic intervals \(I\) of average \(\langle |g|\rangle_I>\lambda\).
They are disjoint, their total length is at most \(\|g\|_1/\lambda\),
and their averages are at most \(2\lambda\) by the parent bound.
Large ancestor averages tend to zero since \(g\in L^1\). Off the union
\(|g|\le\lambda\) almost everywhere by differentiation. Replacing \(g\)
on these intervals by its complex average gives \(g_0\) with
\(\|g_0\|_2^2\le 2\lambda\|g\|_1\) and
\(g=g_0+\sum_I b_I\), where each \(b_I\) has integral zero, is
supported on \(I\), and \(\sum_I\|b_I\|_1\le2\|g\|_1\).
Outside the interval with the same center \(c_I\) and triple length,
subtract \(K(x-c_I)\) inside \(K*b_I(x)\). For \(y\in I\),
\[
 \int_{|x-c_I|>3|I|/2}
 \frac{|y-c_I|}{|x-c_I|^2}\dd x\le1 .
\]
Thus the \(L^1\) norm of \(K*b_I\) there is at most \(\|b_I\|_1\).
The infinite sum agrees in \(L^1\) with convolution of the bad part
since \(K\in L^1\). The union of the tripled intervals, Chebyshev on
the good part in \(L^2\), and Chebyshev outside the union on the bad
part give \(|\{|K*g|>\lambda\}|\le C\|g\|_1/\lambda\).
For \(1<q<2\) split \(g\) at height \(\lambda\). Apply the weak
\(L^1\) bound to \(g\ind_{\{|g|>\lambda\}}\) and the \(L^2\) bound
to the rest, each at output height \(\lambda/2\). Multiply by
\(q\lambda^{q-1}\) and integrate. The two inner integrals after
Tonelli are \(\int_0^{|g|}\lambda^{q-2}\dd\lambda\) and
\(\int_{|g|}^\infty\lambda^{q-3}\dd\lambda\), yielding the desired
moment. For \(q>2\), the reflected conjugate kernel satisfies the
same hypotheses, so duality gives the result.
\end{proof}

\begin{lemma}[Signed sums and square functions]\label{bp:profile-square}
\lean{TwistedParaproduct.profileSquare}
If \(\int a=0\) and \(Q(a)<\infty\), then for any finite
\(I\subset\mathbb Z\), \(|z_k|\le1\), and \(1<q<\infty\),
\[
 \left\|\sum_{k\in I} z_k a_k*_i f\right\|_q+
 \left\|\Big(\sum_{k\in I}|a_k*_i f|^2\Big)^{1/2}\right\|_q
 \le C_q Q(a)\|f\|_q .
\]
\end{lemma}
\begin{proof}
\uses{bp:scalar-cz}
Cancellation and \(|e^{it}-1|\le|t|\) give
\(|\widehat a(\zeta)|\le C Q(a)|\zeta|\). Integration by parts,
using \(a(\pm\infty)=0\) and \(a'\in L^1\), gives the bound
\(C Q(a)/|\zeta|\). The sum of
\(\min(2^{-k}|\zeta|,2^k/|\zeta|)\) over all integers is bounded by
two geometric series. The kernel \(K=\sum_{k\in I}z_k a_k\) therefore
satisfies \(\|\widehat K\|_\infty\le C Q(a)\). Its almost-everywhere
derivative is bounded, away from zero, by
\[
 Q(a)\sum_{k\in\mathbb Z}2^{2k}(1+2^k|x|)^{-4}
 \le C Q(a)|x|^{-2}.
\]
Split the sum at \(2^k |x|=1\) to obtain the last inequality.
Absolute continuity on intervals gives the difference hypothesis of
Lemma~\ref{bp:scalar-cz}, since every point between \(x-u\) and \(x\)
has modulus at least \(|x|/2\) when \(2|u|\le |x|\). This proves the
signed bound on the line and, by Fubini, on each coordinate.

For independent uniform signs \(\varepsilon_k\) and complex scalars
\(w_k\), put \(X=|\sum_k\varepsilon_k w_k|^2\) and \(S=\sum_k|w_k|^2\).
Expansion gives \(\E X=S\) and \(\E X^2\le3 S^2\), since only
indices paired in one of three ways can survive. For \(S>0\),
\(S/2\le\E[X\ind_{\{X>S/2\}}]\le(\E X^2)^{1/2}
\mathbb P(X>S/2)^{1/2}\). Thus \(S^{q/2}\le12\,2^{q/2}\E X^{q/2}\).
Apply this pointwise and integrate the signed estimate.
\end{proof}

\begin{lemma}[An annular correction profile]\label{bp:tent-inverse}
\lean{TwistedParaproduct.tentInverse}
Fix \(0<c<C<\infty\). There is \(\lambda>0\) such that
\(D(z)=\widehat A(\lambda z/2)-\widehat A(\lambda z)\) is nonzero
on an open neighborhood of \(c\le |z|\le C\). For a Schwartz function
\(\psi\) with Fourier support in that annulus, there is a Schwartz
function \(v\) of integral zero with \(D\widehat v=\widehat\psi\) and
\[
 Q(v)\le C_{c,C}\sum_{j=0}^4\|(\widehat\psi)^{(j)}\|_\infty .
\]
\end{lemma}
\begin{proof}
Evenness and compact support yield \(\widehat A(z)=1-
2\pi^2 z^2\int y^2A(y)\dd y+O(z^4)\). The quadratic coefficient is
strictly negative, so \(\widehat A(w/2)-\widehat A(w)>0\) for
sufficiently small nonzero \(w\). Choose \(\lambda\) accordingly, and
choose a smooth cutoff equal to one on the annulus, supported in a
bounded open set where \(D\ne0\), away from zero. Divide this cutoff
by \(D\), extend it by zero, and multiply by \(\widehat\psi\) to
obtain \(\widehat v\). This smooth compactly supported function
vanishes at zero. Fourier inversion bounds \(v,v'\) by the \(L^1\)
norms of \(\widehat v,2\pi iz\widehat v\); four integrations by parts
give the same bounds for \(y^4v,y^4v'\) using the fourth derivatives.
The claimed estimate follows by the product rule on a fixed compact set.
\end{proof}

\begin{lemma}[Averaging translated dyadic grids]\label{bp:random-grid}
\lean{TwistedParaproduct.randomGrid}
For \(u\le k\le v\), translate the grid of intervals of length
\(\lambda_0 2^{-k}\) by a single variable uniform on
\([0,\lambda_0 2^{-u})\), with \(\lambda_0>0\).
The expected conditional expectation at level \(k\) is convolution
with \(L^{-1}A(\cdot/L)\), \(L=\lambda_0 2^{-k}\).
On independently translated coordinate grids, allowing a separate
fixed dilation in each coordinate, the finite dyadic estimate
\ref{bp:dyadic} has its ordinary constant.
\end{lemma}
\begin{proof}
\uses{bp:dyadic}
The translate modulo \(L\) is uniform since the largest length is an
integer multiple of \(L\). The probability that \(x,y\) belong to the
same grid interval is \((1-|x-y|/L)_+\), by measuring the possible
location of the left boundary modulo \(L\). Multiply by \(1/L\) and
integrate against the input. The expected difference of two successive
levels is the difference of the two expected kernels. For the estimate
on the changed grids, use the same affine change of variables in
every input and in the tester. Its determinant cancels because the
reciprocals of all form exponents sum to one. Translation and
coordinate dilation conjugate the conditional expectations exactly.
\end{proof}

\begin{lemma}[Testing the output against Schwartz functions]\label{bp:dual}
\lean{TwistedParaproduct.dual}
For \(1<p<\infty\), \(p'=p/(p-1)\), and \(g\in L^p(X)\),
\[
 \norm g_p=\sup\{|\textstyle\int gh|:
                       h\in\Sch(X;\C),\ \norm h_{p'}\le1\}.
\]
\end{lemma}
\begin{proof}
H\"older gives one inequality. If \(g\ne0\), the function that is
\(\overline g\,|g|^{p-2}/\norm g_p^{p-1}\) on \(g\ne0\) and zero
elsewhere has \(p'\)-norm one and gives equality. Approximate it in
\(L^{p'}\) by Schwartz functions and divide each approximant by the
maximum of one and its norm. H\"older passes the pairing to the limit.
For \(g=0\) both sides vanish.
\end{proof}

\begin{lemma}[Separated finite sums]\label{bp:separated-cone}
\lean{TwistedParaproduct.separatedCone}
Let \(\widehat\psi\) have support in \(c\le |z|\le C\), and let
\(\psi,\phi_i\) be Schwartz profiles on \(\R\). Write
\[
 P(\psi)=Q(\psi)+\sum_{j=0}^4\|(\widehat\psi)^{(j)}\|_\infty .
\]
Fix \(b\in\{1,\ldots,d\}\). For any finite interval \(I\) of integers
and coefficients \(|c_k|\le1\),
\[
 \left\|\sum_{k\in I}c_k(\psi_k*_b f_b)
            \prod_{i\ne b}(\phi_{i,k}*_i f_i)\right\|_p
 \le C_{d,\vec p,c,C}\,P(\psi)\prod_{i\ne b}Q(\phi_i)
                             \prod_{i=1}^d\|f_i\|_{p_i}.
\]
\end{lemma}
\begin{proof}
\uses{bp:random-grid,bp:tent-inverse,bp:profile-square,bp:dual,bp:profile-max,bp:dyadic}
First take the low profiles all equal to \(A\). Use
Lemma~\ref{bp:random-grid} with dilation \(1\) in those coordinates
and \(\lambda\) from Lemma~\ref{bp:tent-inverse} in the distinguished one.
In the finite dyadic estimate replace the coefficients by
\(c_k\varepsilon_k\) and the distinguished input by
\(\sum_{\ell\in I}\varepsilon_\ell v_\ell*_b f_b\). Its
\(L^{p_b}\) norm is bounded uniformly by Lemma~\ref{bp:profile-square}.
In the form paired against a Schwartz tester, expectation over the signs
retains precisely \(\ell=k\). At each such \(k\), independence of the
grid translations factors their expectation over coordinates. The high
factor has multiplier \(D(2^{-k}\zeta)\widehat v(2^{-k}\zeta)
=\widehat\psi(2^{-k}\zeta)\). The low factors have profiles \(A\).
These exchanges are justified by finite sums and Fubini: each grid
average contracts \(L^{p_i}\), each difference has norm at most two,
and H\"older bounds the integral of the absolute value at each index.
Each finite output belongs to \(L^p\) already by the individual
one-coordinate \(L^1\)-kernel estimates and H\"older. Apply
Lemma~\ref{bp:dual} to obtain the low-tent bound.

For general low profiles write \(m_i=\int\phi_i\) and telescope
from \(m_i A\) to \(\phi_i\) in a fixed order. The initial product
costs \(\prod_{i\ne b}|m_i|\le C_d\prod_{i\ne b}Q(\phi_i)\).
An error term has a mean-zero low profile
\(a_i=\phi_i-m_i A\), with \(Q(a_i)\le C Q(\phi_i)\), and the
mean-zero high profile \(\psi\). Bound the sum over \(k\) pointwise by
the product of their two square functions and the scale suprema of
the remaining factors. Lemmas~\ref{bp:profile-max} and
\ref{bp:profile-square} and H\"older give the asserted bound.
\end{proof}

\subsection{Rescaled symbols and Fourier series}
We now return to the symbol and Fourier convention in \eqref{bp:op}.

\begin{lemma}[A finite family of conical cutoffs]\label{bp:cone-cutoffs}
\lean{TwistedParaproduct.coneCutoffs}
There are nonnegative \(\theta,\theta_1,\ldots,\theta_d
\in C_c^\infty(\R^d)\) with
\[
 \theta=\sum_b\theta_b,\qquad
 \sum_{k\in\mathbb Z}\theta(2^{-k}\eta)=1\quad(\eta\ne0),\qquad
 \operatorname{supp}\theta_b\subseteq
 \{1/2\le |z|\le2,\ |z_b|\ge 1/(4\sqrt d)\}.
\]
One may take \(\theta\) radial. For every \(b\) there are
one-dimensional smooth cutoffs \(\gamma_{b,i}\) whose product equals
one on \(\operatorname{supp}\theta_b\), with all supports in \((-3,3)\)
and the support for \(i=b\) in \(\{1/(8\sqrt d)\le |s|\le3\}\).
\end{lemma}
\begin{proof}
Take a radial smooth cutoff \(\chi\), nonincreasing in radius,
equal to one on the unit ball and zero outside radius two. Set
\(\theta(z)=\chi(z)-\chi(2z)\); the sum telescopes. Take smooth
\(w\ge0\) vanishing on \([0,1/(4d)]\) and positive on \([1/d,1]\).
Multiply \(\theta(z)\) by \(w(z_b^2/|z|^2)/\sum_iw(z_i^2/|z|^2)\)
and extend by zero near the origin. At least one squared ratio is
\(\ge1/d\), so the denominator is positive. The supports are as
claimed. For the product cutoffs choose functions equal to one on
\([-2,2]\) in low coordinates and on \(1/(4\sqrt d)\le |s|\le2\)
in the distinguished coordinate, inside the indicated larger sets.
\end{proof}

\begin{lemma}[Uniform separated expansions]\label{bp:mode-expansion}
\lean{TwistedParaproduct.modeExpansion}
Set \(M_{k,b}(z)=\theta_b(z)m(2^kz)\), extended by zero near zero.
For each integer \(s>d\) it admits the expansion on \(\R^d\)
\[
 M_{k,b}(z)=\sum_{n\in\mathbb Z^d} a_{k,b,n}
  \prod_i\bigl(\gamma_{b,i}(z_i)e^{2\pi i n_i z_i/8}\bigr),
 \qquad |a_{k,b,n}|\le C_{d,s}[m]_s(1+|n|)^{-s}.
\]
For fixed \(b\) let \(\widehat{\phi_{n,i}}(z)=
\gamma_{b,i}(z)e^{2\pi i n_i z/8}\). Then
\(Q(\phi_{n,i})\le C_d(1+|n_i|)^4\), and also
\(P(\phi_{n,b})\le C_d(1+|n_b|)^4\).
\end{lemma}
\begin{proof}
\uses{bp:cone-cutoffs}
The chain rule bounds all derivatives of \(M_{k,b}\) through order
\(s\) by \(C_{d,s}[m]_s\) on its fixed annulus. For a Fourier
coefficient on the period cube \([-4,4]^d\) with normalized measure
\(8^{-d}\dd z\), integrate \(s\) times by parts in a coordinate
with \(|n_i|\ge |n|/\sqrt d\) if \(n\ne0\). Boundary terms vanish.
For \(n=0\) use the integral bound. The resulting coefficient
estimate gives absolute uniform convergence. Its sum equals the
periodization, since the continuous difference has all Fourier
coefficients zero. Indeed trigonometric polynomials are uniformly
dense on the period torus by Stone--Weierstrass, so its squared
modulus has zero integral. Multiplying by \(\prod_i\gamma_{b,i}\)
restricts to the original cube and preserves \(M_{k,b}\) there.
The inverse transform of each modulated cutoff is
\(\check\gamma_{b,i}(\cdot+n_i/8)\). Since \(1+|y|\le
(1+|y+n_i/8|)(1+|n_i|/8)\), translation costs the claimed fourth
power in \(Q\). The product rule through four derivatives of the
modulated frequency cutoff gives the additional terms in \(P\).
\end{proof}

\begin{lemma}[Full smooth multiplier estimate]\label{bp:full-multiplier}
\lean{TwistedParaproduct.fullMultiplier}
For \(s=5d+2\), every smooth symbol on \(\R^d\setminus\{0\}\)
of finite \([m]_s\) satisfies on Schwartz inputs
\[
 \|T_m(\vec f)\|_p\le C_{d,\vec p}[m]_s\prod_i\|f_i\|_{p_i}.
\]
\end{lemma}
\begin{proof}
\uses{bp:mode-expansion,bp:separated-cone,bp:schwartz-Fourier,bp:cone-cutoffs}
First use the symbol \(\sum_{k\in I}m(\eta)\theta(2^{-k}\eta)\)
for a finite integer interval \(I\). For fixed \(b,n\), apply
Lemma~\ref{bp:mode-expansion} at \(z=2^{-k}\eta\). Each factor
depends only on the designated coordinate \(\xi_i^i\) in its own
input. Fourier inversion therefore gives the physical factor
\(\phi_{n,i,k}*_i f_i\) with \(\phi_{n,i,k}(y)=
2^k\phi_{n,i}(2^k y)\). Fubini holds since the input transforms are
integrable and the fixed frequency profiles are bounded. After
rescaling \(a_{k,b,n}\) by their common upper bound, the sum over
\(k\) is estimated by Lemma~\ref{bp:separated-cone}, at cost
\[
 C_{d,\vec p}[m]_s(1+|n|)^{4d-s}\prod_i\|f_i\|_{p_i}.
\]
Zero coefficients cost zero. This is summable over \(n\), and \(b\)
has \(d\) possible values. For each finite \(I\) the mode expansion
converges pointwise in the Fourier integral by its absolute
coefficient bound. The triangle inequality for finite mode sums and
Fatou bound the resulting operator uniformly in \(I\).
Now take \(I=[-j,j]\cap\mathbb Z\). The symbols converge to \(m\)
off zero and have modulus at most \(\|m\|_\infty\), since \(\theta\)
is nonnegative. Dominated convergence gives pointwise convergence
in the Fourier integral to the untruncated \(T_m\). Fatou proves
the assertion.
\end{proof}

\subsection{Physical kernels and the principal value}

\begin{lemma}[An integrable physical kernel]\label{bp:physical-integrable}
\lean{TwistedParaproduct.physicalIntegrable}
For \(L\in L^1(\R^d)\), put
\[
 U_L(\vec f)(x)=\int_{\R^d} L(y)\prod_{i=1}^d f_i(x-y_i e_i)\dd y .
\]
Then \(\|U_L(\vec f)\|_p\le\|L\|_1\prod_i\|f_i\|_{p_i}\).
On Schwartz inputs \(U_L=T_{\widehat L}\) pointwise.
\end{lemma}
\begin{proof}
Apply Minkowski in \(y\) and H\"older in \(x\) for each \(y\), using
translation invariance of the input norms. On Schwartz inputs expand
the factors by Fourier inversion. The absolute integral before
exchanging integration is at most \(\|L\|_1\prod_i\|\widehat f_i\|_1\).
The \(y\)-integral gives \(\widehat L(\xi_1^1,\ldots,\xi_d^d)\).
\end{proof}

\begin{lemma}[Off-origin kernel]\label{bp:mihlin-kernel}
\lean{TwistedParaproduct.mihlinKernel}
The tempered distribution \(\mathcal K=\mathcal F^{-1}m\) of a
smooth Mihlin symbol is represented away from zero by a \(C^1\)
function \(K\) satisfying
\[
 |K(y)|+|y|\,|\nabla K(y)|\le C_d[m]_{d+2}|y|^{-d}\quad(y\ne0).
\]
\end{lemma}
\begin{proof}
\uses{bp:cone-cutoffs}
Use Lemma~\ref{bp:cone-cutoffs} and put \(L_k=\mathcal F^{-1}
(m\,\theta(2^{-k}\cdot))\). For \(\ell=0,1\), Fourier inversion
at rescaled frequency and \(d+2\) integrations by parts in a
coordinate of largest \(|y_i|\) give
\[
 |\nabla^\ell L_k(y)|\le C_d[m]_{d+2}
 2^{k(d+\ell)}(1+2^k|y|)^{-d-2}.
\]
For \(2^k|y|\le1\) use the direct integral bound. Rescaled
derivatives of the amplitude, including the linear frequency factor
for \(\ell=1\), have bounded \(L^1\) norms through order \(d+2\)
on a fixed annulus by the chain and product rules. Splitting the
geometric sum at \(2^k|y|=1\) gives \(C_d[m]_{d+2}|y|^{-d-\ell}\).
The series and its first derivatives converge locally uniformly off
zero. Explicitly, \(\langle\mathcal K,g\rangle=\int m(\eta)
\widehat g(-\eta)\dd\eta\) for Schwartz \(g\), with bilinear pairing.
For smooth \(g\) compactly supported away from zero, integrals of
the partial kernel sums against \(g\) converge to \(\int Kg\) by
local uniform convergence. Fourier inversion of every frequency
piece and dominated convergence against \(\widehat g(-\cdot)\)
identify the same limit with \(\langle\mathcal K,g\rangle\).
\end{proof}

\begin{lemma}[Comparison with mollified kernels]\label{bp:hard-comparison}
\lean{TwistedParaproduct.hardComparison}
Fix \(\rho\in C_c^\infty(\R^d)\) of integral one and supported in the
unit ball; put \(\rho_r(y)=r^{-d}\rho(y/r)\). For \(K,\mathcal K\)
as in Lemma~\ref{bp:mihlin-kernel},
\[
 \left\|\ind_{\{|y|>r\}}K-\mathcal K*\rho_r\right\|_1
          \le C_{d,\rho}[m]_{d+2}\qquad(r>0).
\]
\end{lemma}
\begin{proof}
\uses{bp:mihlin-kernel}
Apply this explicit pairing to \(g(z)=\rho_r(y-z)\):
\(\widehat g(-\eta)=e^{2\pi i y\cdot\eta}\widehat\rho(r\eta)\).
The mollification is consequently the ordinary inverse transform of
\(m(\eta)\widehat\rho(r\eta)\), and its modulus is at most
\(r^{-d}\|m\|_\infty\|\widehat\rho\|_1\). This bounds the integral
on \(|y|\le2r\) uniformly. For \(K\) on \(r<|y|\le2r\) integrate
the size bound. On \(|y|>2r\) the distributional mollification is
the ordinary local integral of \(K\), since its test function
has support away from the origin. Subtract \(K(y)\int\rho_r\).
The gradient bound along each segment gives
\(C_{d,\rho}[m]_{d+2}r|y|^{-d-1}\), whose integral on this exterior
region is uniformly bounded.
\end{proof}

\begin{lemma}[Uniform hard annuli]\label{bp:hard-general}
\lean{TwistedParaproduct.hardGeneral}
For a smooth Mihlin symbol \(m\) and its off-origin kernel \(K\),
\[
 \sup_{0<\epsilon<R<\infty}
 \left\|U_{\ind_{\{\epsilon<|y|<R\}}K}(\vec f)\right\|_p
 \le C_{d,\vec p}[m]_{5d+2}\prod_i\|f_i\|_{p_i}.
\]
\end{lemma}
\begin{proof}
\uses{bp:hard-comparison,bp:full-multiplier,bp:physical-integrable}
Subtract the discrepancies of Lemma~\ref{bp:hard-comparison} at
\(r=\epsilon,R\). Up to a uniformly controlled \(L^1\) error the
hard annulus is \(\mathcal K*\rho_\epsilon-\mathcal K*\rho_R\).
This difference is itself integrable and has Fourier transform
\(m(\eta)(\widehat\rho(\epsilon\eta)-\widehat\rho(R\eta))\)
almost everywhere by distributional Fourier inversion. Both sides
are continuous away from zero, so equality holds there pointwise:
every nonempty open ball has positive measure. For each
\(\alpha\), the quantities \(|\eta|^{|\alpha|}
|\partial^\alpha[\widehat\rho(r\eta)]|\) are uniformly bounded in
\(r>0,\eta\ne0\) by the chain rule and rapid decay. The product rule
bounds the degree-\(s\) seminorm of this symbol by
\(C_{d,s,\rho}[m]_s\). Apply Lemma~\ref{bp:full-multiplier} and
Lemma~\ref{bp:physical-integrable} to the difference and the
integrable error. Fix one mollifier in every dimension.
\end{proof}

\begin{lemma}[Multiplier of a homogeneous principal value]
\label{bp:homogeneous-symbol}
\lean{TwistedParaproduct.homogeneousSymbol}
Let \(\Omega\in C^\infty(S^{d-1})\) have surface integral zero, and
put \(K(y)=|y|^{-d}\Omega(y/|y|)\) off zero. Its unmodified spherical
principal-value distribution \(D_\Omega\), defined by
\[
 \langle D_\Omega,g\rangle=
 \int K(y)\big(g(y)-g(0)\ind_{\{|y|\le1\}}\big)\dd y
 \quad(g\ \text{Schwartz}),
\]
has a smooth Mihlin symbol \(m_\Omega=\widehat{D_\Omega}\) off zero
and is the inverse transform of this bounded multiplier on all of
\(\R^d\) in tempered distributions.
\end{lemma}
\begin{proof}
\uses{bp:cone-cutoffs}
The defining integral converges by the estimate \(g(y)-g(0)=O_g(|y|)\)
near zero and rapid decrease at infinity. These estimates in terms
of Schwartz seminorms give continuity on the Schwartz space.
Cancellation on spheres identifies it with hard spherical principal
values. Use the radial annulus \(\theta\) of Lemma~\ref{bp:cone-cutoffs}
in the physical variable and put \(k(y)=K(y)\theta(y)\), smoothly
extended by zero. It is smooth, compactly supported, and has integral
zero by polar coordinates. Write \(k_r(y)=r^{-d}k(y/r)\), so
\(k_r(y)=K(y)\theta(y/r)\). The sums \(\sum_{|\ell|\le j}k_{2^\ell}\)
converge to \(D_\Omega\) on Schwartz tests: subtract
\(g(0)\ind_{\{|y|\le1\}}\) in each annulus by radial cancellation
and use \(0\le\sum_{|\ell|\le j}\theta(2^{-\ell}y)\le1\) for domination.

Off zero the series \(m_\Omega(\eta)=\sum_{\ell\in\mathbb Z}
\widehat k(2^\ell\eta)\) and all its derivative series converge
locally uniformly. With \(r=2^\ell\) and \(h=|\alpha|\), for
\(r|\eta|\le1\) the summand has modulus at most \(C r|\eta|\) if
\(h=0\) (subtract one inside the transform), and its derivative of
order \(h\ge1\) has modulus at most \(C_h r^h\) by differentiating
the Fourier integral. For \(r|\eta|>1\), integration by parts \(h+2\)
times in a coordinate of largest \(|\eta_i|\) bounds that derivative
by \(C_h r^h(r|\eta|)^{-h-2}\). These constants use only finitely
many derivatives of \(k\) for each \(h\). Summing the two geometric
series gives \(C_h|\eta|^{-h}\), uniformly for finite partial sums
as well. In particular the frequency sums for \(h=0\) are uniformly
bounded globally. Dominated convergence against Schwartz tests now
shows \(\widehat{D_\Omega}=m_\Omega\) as distributions on the whole
space. This also fixes the additive constant in the multiplier
through the specified spherical normalization.
\end{proof}

\begin{lemma}[The untruncated principal value on Schwartz inputs]
\label{bp:hard-homogeneous}
\lean{TwistedParaproduct.hardHomogeneous}
For \(\Omega,K,m_\Omega\) as in Lemma~\ref{bp:homogeneous-symbol}
and Schwartz inputs, the limit
\[
 \operatorname{p.v.}\int_{\R^d}
  K(y)\prod_i f_i(x-y_i e_i)\dd y
 =\lim_{\epsilon\downarrow0,\ R\uparrow\infty}
     U_{\ind_{\{\epsilon<|y|<R\}}K}(\vec f)(x)
 =T_{m_\Omega}(\vec f)(x)
\]
holds for every \(x\). The hard annuli converge also in \(L^p\);
the principal value satisfies
\(\|T_{m_\Omega}(\vec f)\|_p\le
C_{d,\vec p,\Omega}\prod_i\|f_i\|_{p_i}\).
\end{lemma}
\begin{proof}
\uses{bp:homogeneous-symbol,bp:physical-integrable,bp:full-multiplier,bp:hard-general,bp:schwartz-Fourier}
Put \(G_x(y)=\prod_i f_i(x-y_i e_i)\). For every fixed \(x\),
\(G_x\) is Schwartz in \(y\): each factor is a Schwartz function
of its own \(y_i\) at the indicated fixed spectator coordinates.
By spherical cancellation the truncated integral, for
\(\epsilon<1<R\), is the same with
\(G_x(y)-G_x(0)\ind_{\{|y|\le1\}}\) in place of \(G_x(y)\).
Absolute convergence gives the pointwise limit
\(\langle D_\Omega,G_x\rangle\) along all pairs of endpoints.
For the finite physical sums \(L_j=\sum_{|\ell|\le j}k_{2^\ell}\)
of Lemma~\ref{bp:homogeneous-symbol}, Fubini and Fourier inversion
give \(U_{L_j}(\vec f)(x)=T_{\widehat L_j}(\vec f)(x)\) by
Lemma~\ref{bp:physical-integrable}. The left side tends to
\(\langle D_\Omega,G_x\rangle\) by the distributional limit.
On the right the symbols are uniformly bounded and approach
\(m_\Omega\) off zero. The product of the input Fourier transforms
is integrable, so dominated convergence identifies the right limit
with \(T_{m_\Omega}(\vec f)(x)\). The norm bound now follows from
Lemma~\ref{bp:full-multiplier}; Lemma~\ref{bp:hard-general} also
gives the uniform bound for the original hard annuli.

For norm convergence, when \(|y|\le1\) the product rule and the
mean value theorem give
\(|G_x(y)-G_x(0)|\le C_L |y|(1+|x|)^{-L}\) for any \(L>0\).
For large \(y\) put \(B=\max_i|x-y_i e_i|\). For each coordinate
\(j\), choose \(i\ne j\) to see \(|x_j|\le B\); then
\(|y_i|\le|x_i|+|(x-y_i e_i)_i|\le2B\).
The largest argument therefore controls \(|x|+|y|\) up to a
dimension constant. Bounding the corresponding Schwartz factor by
rapid decrease and all others by their suprema gives
\[
 |G_x(y)|\le C_L(1+|x|)^{-L}(1+|y|)^{-1}
 \qquad (|y|\ge1).
\]
By polar integration the discarded regions in the subtracted integral
are bounded in modulus by \(C_{L,\Omega}(\epsilon+R^{-1})(1+|x|)^{-L}\).
Taking \(Lp>d\) gives convergence in \(L^p\).
\end{proof}

\begin{proof}[Proof of Theorem~\ref{bp:main}]
\uses{bp:full-multiplier,bp:frechet,bp:homogeneous-symbol,bp:hard-homogeneous}
Use Lemma~\ref{bp:full-multiplier} with \(s=5d+2\) in
\eqref{bp:op}. Lemma~\ref{bp:frechet} identifies its quantitative
hypothesis with the bound in the Lean target. For a homogeneous
mean-zero kernel, Lemma~\ref{bp:homogeneous-symbol} supplies a symbol
with finite \([m_\Omega]_s\), and Lemma~\ref{bp:hard-homogeneous}
identifies the untruncated principal value pointwise with
\(T_{m_\Omega}\) on the same Schwartz inputs. Substitution gives the
second inequality. All constants in the argument are independent of
the scale endpoints.
\end{proof}

\clearpage
\appendix
\section{Lean statement of the Fourier estimate}\label{bp:lean-target}
The following code elaborates using the library revisions in
Appendix~\ref{sec:tpb-source-map}. It defines the proposition that the
formalization must prove. The bound uses an explicit nonnegative upper
bound \(B\) because real suprema in Lean are total even for unbounded
sets. Lemma~\ref{bp:frechet} gives the conversion. The value of the
symbol at zero can be arbitrary by Lemma~\ref{bp:schwartz-Fourier}.
\begin{lstlisting}
import Mathlib.Analysis.Distribution.SchwartzSpace.Fourier
import Mathlib.MeasureTheory.Function.LpSeminorm.Defs

open MeasureTheory
open scoped BigOperators FourierTransform

namespace TwistedParaproduct
abbrev E (d : ℕ) := EuclideanSpace ℝ (Fin d)

noncomputable def twisted {d : ℕ} (m : E d → ℂ)
    (f : Fin d → SchwartzMap (E d) ℂ) (x : E d) : ℂ :=
  ∫ ξ : Fin d → E d,
    m (WithLp.toLp 2 (fun j => ξ j j)) *
      (∏ j, (𝓕 (f j)) (ξ j)) *
      Complex.exp (Complex.I * Complex.ofReal
        (2 * Real.pi * inner ℝ x (∑ j, ξ j)))

def BanachTarget : Prop :=
  ∀ (d : ℕ), 2 ≤ d → ∀ (p : Fin d → ℝ) (r : ℝ),
    (∀ j, 1 < p j) → 1 < r → r⁻¹ = ∑ j, (p j)⁻¹ →
    ∃ (N : ℕ) (C : ℝ), 0 < C ∧ ∀ (m : E d → ℂ) (B : ℝ),
      0 ≤ B → (∀ n : ℕ, ContDiffOn ℝ n m {x | x ≠ 0}) →
      (∀ k ≤ N, ∀ z : E d, z ≠ 0 →
        ‖z‖ ^ k * ‖iteratedFDeriv ℝ k m z‖ ≤ B) →
      ∀ f : Fin d → SchwartzMap (E d) ℂ,
        eLpNorm (twisted m f) (ENNReal.ofReal r) volume ≤
          ENNReal.ofReal (C * B) *
            ∏ j, eLpNorm (fun x => f j x) (ENNReal.ofReal (p j)) volume

#check TwistedParaproduct.BanachTarget
end TwistedParaproduct
\end{lstlisting}
For the spatial formulation, represent the untruncated value by
\(\int K(y)(G_x(y)-G_x(0)\ind_{\{|y|\le1\}})\dd y\), where
\(G_x(y)=\prod_i f_i(x-y_ie_i)\). Lemma~\ref{bp:hard-homogeneous}
proves absolute integrability and equality with the principal value.
The corresponding two-endpoint limit in Lean uses the product of the
right-neighborhood filter at zero and \code{Filter.atTop} on the outer
radius. The equation is a pointwise equality for each \(x\).


\section{Lean source map}\label{sec:tpb-source-map}
\begingroup\raggedright
The references below use Lean \nolinkurl{v4.35.0-rc3}, \href{https://github.com/roos-j/lean-spherical/blob/800b71afca91e42809c21e6e2fd9ae4f82467082/lakefile.toml#L1}{the lean-spherical package at} \texttt{800b71afca91e42809c21e6e2fd9ae4f82467082}, and \href{https://github.com/roos-j/lean-spherical/blob/800b71afca91e42809c21e6e2fd9ae4f82467082/lake-manifest.json#L1}{its resolved mathlib dependency at} \texttt{c55e6e786f49471c72fbddbec5415808896aec1e}. The links select the precise source files and lines at these commits. The remaining interfaces described below are formalization work for the twisted operator.

\paragraph{Coordinates, products, and density.}
\nolinkurl{PiLp.volume_preserving_ofLp} and \nolinkurl{PiLp.volume_preserving_toLp} identify Lebesgue measure in Euclidean coordinates with finite product volume. Use these identifications, a permutation placing the active coordinate first, and product integration before applying a one-dimensional estimate on each fiber. \nolinkurl{MeasureTheory.integral_prod} applies to an integrable function; its nonnegative counterpart is \nolinkurl{MeasureTheory.lintegral_prod}. \nolinkurl{ENNReal.lintegral_prod_norm_pow_le} handles a finite family of weighted factors. \nolinkurl{SchwartzMap.denseRange_toLpCLM} supplies the Schwartz approximation in Lemma~\ref{bp:dual} at the tester exponent. Simple-function density is used separately in Proposition~\ref{bp:dyadic}. \href{https://github.com/leanprover-community/mathlib4/blob/c55e6e786f49471c72fbddbec5415808896aec1e/Mathlib/MeasureTheory/Measure/Haar/InnerProductSpace.lean#L134}{Coordinates}, \href{https://github.com/leanprover-community/mathlib4/blob/c55e6e786f49471c72fbddbec5415808896aec1e/Mathlib/MeasureTheory/Integral/Prod.lean#L440}{Fubini}, \href{https://github.com/leanprover-community/mathlib4/blob/c55e6e786f49471c72fbddbec5415808896aec1e/Mathlib/MeasureTheory/Measure/Prod.lean#L1009}{Tonelli}, \href{https://github.com/leanprover-community/mathlib4/blob/c55e6e786f49471c72fbddbec5415808896aec1e/Mathlib/MeasureTheory/Integral/MeanInequalities.lean#L241}{finite H\"older inequality}, \href{https://github.com/leanprover-community/mathlib4/blob/c55e6e786f49471c72fbddbec5415808896aec1e/Mathlib/Analysis/Distribution/SchwartzSpace/Basic.lean#L1377}{Schwartz density}.

\paragraph{Martingales.}
In namespace \nolinkurl{MeasureTheory}, \nolinkurl{martingale_condExp} gives the conditional-expectation process. \nolinkurl{Adapted.isStoppingTime_hittingBtwn_isStoppingTime} allows a bounded stopping time as the starting index for hitting a fixed measurable set. Package the displacement from the last held value as an adapted process and handle the terminal index and reset convention separately. \nolinkurl{Martingale.stoppedValue_ae_eq_condExp_of_le} gives the almost-everywhere conditional-expectation identity for ordered bounded stopping times with countable time index and the stated sigma-finiteness hypothesis on the trimmed measure. \nolinkurl{maximal_ineq} is a level-set inequality on a finite measure space for a nonnegative submartingale. For coordinate Lebesgue measure one can first work on coarsest finite intervals and then sum over them. The strong maximal moments, square-function moments, signed transforms and reset estimates used by the stopping argument need the scalar estimates developed in this blueprint. \href{https://github.com/leanprover-community/mathlib4/blob/c55e6e786f49471c72fbddbec5415808896aec1e/Mathlib/Probability/Martingale/Basic.lean#L141}{Conditional expectations}, \href{https://github.com/leanprover-community/mathlib4/blob/c55e6e786f49471c72fbddbec5415808896aec1e/Mathlib/Probability/Process/HittingTime.lean#L433}{hitting times}, \href{https://github.com/leanprover-community/mathlib4/blob/c55e6e786f49471c72fbddbec5415808896aec1e/Mathlib/Probability/Martingale/OptionalSampling.lean#L141}{optional sampling}, \href{https://github.com/leanprover-community/mathlib4/blob/c55e6e786f49471c72fbddbec5415808896aec1e/Mathlib/Probability/Martingale/OptionalStopping.lean#L144}{maximal inequality}.

\paragraph{Maximal functions and finite signs.}
In namespace \nolinkurl{Auto.HardyLittlewoodMaximal}, \nolinkurl{dyadic_ball_maximal_raw_weak_one} concerns the extended nonnegative supremum over \emph{centered balls of dyadic radii}. \nolinkurl{dyadic_hardy_littlewood_maximal_strong_type_schwartz} controls the actual pointwise maximal function on Schwartz input. \nolinkurl{dyadic_hardy_littlewood_maximal_Lp} is the bounded extension on \(L^p\) classes. For a dyadic grid interval containing the evaluation point, compare its average to the centered ball of radius the interval length. For convolution with a low profile, sum the centered-ball estimates on its decay annuli. \nolinkurl{Auto.LittlewoodPaley.rademacher_lower_moment} is a finite-powerset inequality for complex coefficients at every real moment order \(p\ge0\); divide by the number of subsets to obtain the uniform-sign expectation. It applies to the signed low-profile error as well as to smooth frequency bands. \href{https://github.com/roos-j/lean-spherical/blob/800b71afca91e42809c21e6e2fd9ae4f82467082/LeanSpherical/Auto/HardyLittlewoodMaximal.lean#L330}{Weak estimate}, \href{https://github.com/roos-j/lean-spherical/blob/800b71afca91e42809c21e6e2fd9ae4f82467082/LeanSpherical/Auto/HardyLittlewoodMaximal.lean#L1043}{Schwartz strong estimate}, \href{https://github.com/roos-j/lean-spherical/blob/800b71afca91e42809c21e6e2fd9ae4f82467082/LeanSpherical/Auto/HardyLittlewoodMaximal.lean#L1326}{\(L^p\) extension}, \href{https://github.com/roos-j/lean-spherical/blob/800b71afca91e42809c21e6e2fd9ae4f82467082/LeanSpherical/Auto/LittlewoodPaley.lean#L1357}{finite-sign lower moment}.

\paragraph{Linear multipliers and singular kernels.}
The scalar theorem \nolinkurl{Auto.MikhlinHormander.mikhlin} takes an indexed family satisfying a \emph{common} \nolinkurl{MikhlinCondition} of order \(N>d\), a finite real exponent \(p>1\), and complex Schwartz data. It returns a finite \nolinkurl{ENNReal} bound uniform in the index. The theorem is exported by \nolinkurl{LeanSpherical.Auto.CalderonZygmund}. For the band-quotient multiplier in one dimension, the indices can include both the signs and the finite interval of scales. Verify the derivatives of the compactly supported quotient and its finite annular overlap. The definition \nolinkurl{MikhlinCondition} includes differentiability off the origin, and \nolinkurl{coordinateMikhlinCondition_to_mikhlinCondition} converts coordinate estimates. In the same namespace, \nolinkurl{pow_norm_fourierInv_le_compactFourierIBPConstant} bounds inverse transforms of compact smooth pieces. \href{https://github.com/roos-j/lean-spherical/blob/800b71afca91e42809c21e6e2fd9ae4f82467082/LeanSpherical/Auto/CalderonZygmund.lean#L5192}{Scalar multiplier theorem}, \href{https://github.com/roos-j/lean-spherical/blob/800b71afca91e42809c21e6e2fd9ae4f82467082/LeanSpherical/Auto/MikhlinHormander.lean#L71}{symbol condition}, \href{https://github.com/roos-j/lean-spherical/blob/800b71afca91e42809c21e6e2fd9ae4f82467082/LeanSpherical/Auto/MikhlinHormander.lean#L152}{coordinate conversion}, \href{https://github.com/roos-j/lean-spherical/blob/800b71afca91e42809c21e6e2fd9ae4f82467082/LeanSpherical/Auto/MikhlinHormander.lean#L581}{Fourier decay}.

For the tent comparison as written in the paper, the finite signed physical kernels are continuous. The applicable entry points in \nolinkurl{Auto.MikhlinHormander} are \nolinkurl{weak_one_finiteConvolutionKernel_of_L2_and_CZTail} and \nolinkurl{exists_uniform_strongLp_rawCore_of_weak_one_square_adjoint}. The first asks for boundedness of each kernel, a quantitative \(L^2\) estimate and its particular bad-atom tail condition. The second takes indexed forward and adjoint operators on Schwartz data, measurability, additivity, the indicated membership and endpoint bounds, and the adjoint pairing. Supply those interfaces for the finite signed kernels and their conjugate reflections. The powerset inequality then turns the signed estimate into a square estimate. The cutoff-specific \nolinkurl{littlewoodPaley_of_mikhlin} concerns the library's prescribed dyadic projections. \href{https://github.com/roos-j/lean-spherical/blob/800b71afca91e42809c21e6e2fd9ae4f82467082/LeanSpherical/Auto/CalderonZygmund.lean#L144}{Kernel \(L^2\) condition}, \href{https://github.com/roos-j/lean-spherical/blob/800b71afca91e42809c21e6e2fd9ae4f82467082/LeanSpherical/Auto/CalderonZygmund.lean#L504}{bad-atom tail condition}, \href{https://github.com/roos-j/lean-spherical/blob/800b71afca91e42809c21e6e2fd9ae4f82467082/LeanSpherical/Auto/CalderonZygmund.lean#L4016}{weak endpoint}, \href{https://github.com/roos-j/lean-spherical/blob/800b71afca91e42809c21e6e2fd9ae4f82467082/LeanSpherical/Auto/CalderonZygmund.lean#L4881}{parameter-uniform adjoint argument}, \href{https://github.com/roos-j/lean-spherical/blob/800b71afca91e42809c21e6e2fd9ae4f82467082/LeanSpherical/Auto/CalderonZygmund.lean#L5228}{Littlewood--Paley corollary}.

\paragraph{Fourier series and the remaining interfaces.}
\nolinkurl{UnitAddTorus.mFourierCoeff_eq_integral} computes coefficients on a translated unit cube, and \nolinkurl{UnitAddTorus.hasSum_mFourier_series_apply_of_summable} recovers a continuous function pointwise once the coefficient family is summable. The lattice estimate \nolinkurl{Auto.Spherical.MSS.mss_integer_lattice_decay_periodization_uniform} takes \(d\) and \(N>d\), with a constant depending on both and uniform in the translate. Identify its lattice subtype with integer coordinate vectors. The remaining cone construction includes the smooth partition, rescaling to the torus, uniform coefficient decay and the translated-profile seminorm estimates. The shifted-grid tent identity and the finite-sign diagonal extraction also need explicit proofs for the chosen coordinate model. \href{https://github.com/leanprover-community/mathlib4/blob/c55e6e786f49471c72fbddbec5415808896aec1e/Mathlib/Analysis/Fourier/AddCircleMulti.lean#L251}{Coefficient integral}, \href{https://github.com/leanprover-community/mathlib4/blob/c55e6e786f49471c72fbddbec5415808896aec1e/Mathlib/Analysis/Fourier/AddCircleMulti.lean#L329}{pointwise reconstruction}, \href{https://github.com/roos-j/lean-spherical/blob/800b71afca91e42809c21e6e2fd9ae4f82467082/LeanSpherical/Auto/Spherical/MSS/FourierCubeLattice.lean#L1311}{lattice summability}.

The new mathematical development consists of the masked box recursion, its drift and Haar bounds, the small-support reduction, and the passage from neighboring restricted form estimates to the multilinear strong bound. The library's \nolinkurl{Auto.MarcinkiewiczInterpolation.hasStrongType_of_hasWeakType} has a single input function and full weak-type hypotheses. The paper's last passage also needs the uniform hard-annulus comparison for the inverse Fourier distribution and the norm convergence argument for homogeneous kernels. \href{https://github.com/roos-j/lean-spherical/blob/800b71afca91e42809c21e6e2fd9ae4f82467082/LeanSpherical/Auto/MarcinkiewiczInterpolation.lean#L3160}{Single-operator interpolation}.
\par\paragraph{Polar integration and Fourier normalization.}
For angular mean zero use unnormalized \nolinkurl{volume.toSphere}.
The polar measure equivalence
\nolinkurl{MeasureTheory.Measure.measurePreserving_homeomorphUnitSphereProd}
splits the punctured space into the sphere and the positive radius
with weight \(r^{d-1}\). Thus radial pieces of a homogeneous kernel
with mean-zero angular factor have integral zero. In Lean one may
express angular smoothness by smoothness of
\(y\mapsto\Omega(y/|y|)\) on the punctured space: restricting to the
sphere recovers the original function, and the radial projection is
smooth there. With the Fourier convention in \eqref{bp:op}, inverse
Fourier distribution pairings use \(\widehat g(-\eta)\) on a test \(g\).
\href{https://github.com/leanprover-community/mathlib4/blob/c55e6e786f49471c72fbddbec5415808896aec1e/Mathlib/MeasureTheory/Constructions/HaarToSphere.lean}{Polar measure equivalence},
\href{https://github.com/leanprover-community/mathlib4/blob/c55e6e786f49471c72fbddbec5415808896aec1e/Mathlib/Analysis/Distribution/SchwartzSpace/Fourier.lean}{Schwartz Fourier transform}.

\endgroup

\end{document}
